% Problems, Chapter 4: Tensors and Differential Forms
% Arfken, Weber & Harris, Mathematical Methods for Physicists, 7th ed.
% Requires: amsmath, amssymb, tikz

\section*{Problems}

\Needspace{9\baselineskip}
\subsection*{4.1 Tensor Analysis}

\Needspace{5\baselineskip}
\par\bigskip
\noindent\textbf{4.1.1}\quad Show that if all the components of any tensor of
any rank vanish in one particular coordinate system, they vanish in all
coordinate systems.

\emph{Note.}~This point takes on special importance in the four-dimensional
(4-D) curved space of general relativity. If a quantity, expressed as a
tensor, exists in one coordinate system, it exists in all coordinate systems
and is not just a consequence of a choice of a coordinate system (as are
centrifugal and Coriolis forces in Newtonian mechanics).
\solution
Let $T$ have $p$ contravariant indices and $q$ covariant indices,
so its rank is $p+q$. At a given point, the components in two coordinate
systems related by a nonsingular change of coordinates obey the tensor
transformation law
\[
T'{}^{i_1\cdots i_p}{}_{j_1\cdots j_q}
=\frac{\partial x'^{i_1}}{\partial x^{a_1}}\cdots
 \frac{\partial x'^{i_p}}{\partial x^{a_p}}
 \frac{\partial x^{b_1}}{\partial x'^{j_1}}\cdots
 \frac{\partial x^{b_q}}{\partial x'^{j_q}}
 T^{a_1\cdots a_p}{}_{b_1\cdots b_q}.
\]
Every repeated index is summed. If every unprimed component on the right
is zero, each term in every sum is zero. Thus every primed component is
zero, regardless of the values of the finite Jacobian factors. For rank
zero the transformation law is simply $T'=T$, with the same conclusion.
This proves the assertion pointwise; if the tensor vanishes throughout a
region, apply the argument at each point of that region.

\Needspace{5\baselineskip}
\par\bigskip
\noindent\textbf{4.1.2}\quad The components of tensor $\mathsf{A}$ are equal
to the corresponding components of tensor $\mathsf{B}$ in one particular
coordinate system, denoted by the superscript $0$; that is,
\[
A_{ij}^{0}=B_{ij}^{0}.
\]
Show that tensor $\mathsf{A}$ is equal to tensor $\mathsf{B}$,
$A_{ij}=B_{ij}$, in all coordinate systems.
\solution
The difference $C_{ij}=A_{ij}-B_{ij}$ is a tensor because the
transformation law is linear. Denote the coordinates of the specified
system by $x_0^a$; the label $0$ (a superscript on components, a
subscript on coordinates) names that system and is not a tensor index.
There $C^0_{ab}=A^0_{ab}-B^0_{ab}=0$. In any other system,
\[
C_{ij}
=\frac{\partial x_0^a}{\partial x^i}
 \frac{\partial x_0^b}{\partial x^j}C^0_{ab}=0.
\]
Consequently $A_{ij}-B_{ij}=0$ in every system.

\Needspace{5\baselineskip}
\par\bigskip
\noindent\textbf{4.1.3}\quad The last three components of a 4-D vector vanish
in each of two reference frames. If the second reference frame is not merely a
rotation of the first about the $x_{0}$ axis, meaning that at least one of the
coefficients $\partial(x')^{i}/\partial x^{0}$ ($i=1,2,3$) is nonzero, show
that the zeroth component vanishes in all reference frames. Translated into
relativistic mechanics, this means that if momentum is conserved in two
Lorentz frames, then energy is conserved in all Lorentz frames.
\solution
Write the vector in the first frame as
$V^\mu=(V^0,0,0,0)$. Its spatial components in the second frame are
\[
V'^i=\frac{\partial x'^i}{\partial x^\mu}V^\mu
=\frac{\partial x'^i}{\partial x^0}V^0,
\qquad i=1,2,3.
\]
By hypothesis all three left sides vanish, and at least one coefficient
$\partial x'^i/\partial x^0$ is nonzero. The corresponding equation
therefore gives $V^0=0$. All four components vanish in the first frame,
so Exercise 4.1.1 gives $V'^\mu=0$ in every frame.

For the conservation statement, apply this argument to the difference
between total final and initial four-momenta,
\[
\Delta P^\mu=(\Delta E/c,\Delta p_x,\Delta p_y,\Delta p_z).
\]
Conservation of three-momentum in the two frames supplies the two sets of
three zero components. Two frames in relative motion satisfy the
hypothesis on the coefficients. Indeed, $\partial x'^i/\partial x^0$,
taken at fixed $x^1,x^2,x^3$, is the rate at which the primed spatial
coordinates of a point at rest in the first frame change with $x^0$; if
all three vanished, every point at rest in the first frame would also be
at rest in the second. For example, for a boost with velocity $\beta c$ along
$x^1$, $x'^1=\gamma(x^1-\beta x^0)$ with $\gamma=(1-\beta^2)^{-1/2}$, so
$\partial x'^1/\partial x^0=-\gamma\beta\ne0$. The argument above then
forces $\Delta E=0$, and hence conservation of all four components in
every Lorentz frame.

\Needspace{5\baselineskip}
\par\bigskip
\noindent\textbf{4.1.4}\quad From an analysis of the behavior of a general
second-rank tensor under $90^{\circ}$ and $180^{\circ}$ rotations about the
coordinate axes, show that an isotropic second-rank tensor in 3-D space must
be a multiple of $\delta_{j}^{i}$.
\solution
In orthonormal Cartesian coordinates, write the components as a
$3\times3$ matrix $T$. A rotation $Q$ gives
$T'_{ij}=Q_{ik}Q_{jl}T_{kl}$, that is, $T'=QTQ^{\mathsf T}$;
isotropy means $QTQ^{\mathsf T}=T$ for every proper rotation.

First use the $180^\circ$ rotations
\[
R_1=\operatorname{diag}(1,-1,-1),\qquad
R_2=\operatorname{diag}(-1,1,-1).
\]
If $R=\operatorname{diag}(s_1,s_2,s_3)$, then
$(RTR^{\mathsf T})_{ij}=s_i s_jT_{ij}$. For $R_1$,
$s_1s_2=s_1s_3=-1$ and $s_2s_3=1$, so $R_1$ changes the signs of
$T_{12},T_{13},T_{21},T_{31}$, and invariance forces each to equal its
negative, hence to vanish. For $R_2$, $s_2s_3=-1$, so $R_2$ likewise
forces $T_{23}=T_{32}=0$.
We have proved, without assuming symmetry in advance, that
$T=\operatorname{diag}(a,b,c)$.

Now use $90^\circ$ rotations about the third and first axes:
\[
Q_3=\begin{pmatrix}0&-1&0\\1&0&0\\0&0&1\end{pmatrix},\qquad
Q_1=\begin{pmatrix}1&0&0\\0&0&-1\\0&1&0\end{pmatrix}.
\]
Multiplying first by $Q$ on the left,
\[
Q_3T=\begin{pmatrix}0&-b&0\\a&0&0\\0&0&c\end{pmatrix},\qquad
Q_1T=\begin{pmatrix}a&0&0\\0&0&-c\\0&b&0\end{pmatrix},
\]
and then by $Q^{\mathsf T}$ on the right gives
\[
Q_3TQ_3^{\mathsf T}=\operatorname{diag}(b,a,c),\qquad
Q_1TQ_1^{\mathsf T}=\operatorname{diag}(a,c,b).
\]
Equality with $T$ requires $a=b$ and $b=c$. Therefore
\[
\boxed{T^i{}_j=\lambda\delta^i_j}.
\]
Conversely $Q(\lambda I)Q^{\mathsf T}=\lambda QQ^{\mathsf T}=\lambda I$,
so every tensor of this form is isotropic.

\Needspace{5\baselineskip}
\par\bigskip
\noindent\textbf{4.1.5}\quad The 4-D fourth-rank Riemann-Christoffel curvature
tensor of general relativity, $R_{iklm}$, satisfies the symmetry relations
\[
R_{iklm}=-R_{ikml}=-R_{kilm}.
\]
With the indices running from $0$ to $3$, show that the number of independent
components is reduced from $256$ to $36$ and that the condition
\[
R_{iklm}=R_{lmik}
\]
further reduces the number of independent components to $21$. Finally, if the
components satisfy an identity $R_{iklm}+R_{ilmk}+R_{imkl}=0$, show that the
number of independent components is reduced to $20$.

\emph{Note.}~The final three-term identity furnishes new information only if
all four indices are different.
\solution
There are initially $4^4=256$ components. Antisymmetry in
$(i,k)$ makes every component with $i=k$ zero and determines a component
with $i>k$ as the negative of one with $i<k$. The six possible ordered
representatives of a pair are
\[
01,\quad02,\quad03,\quad12,\quad13,\quad23.
\]
The same argument applies independently to $(l,m)$. Thus the tensor is
specified by a $6\times6$ array $C_{(ik),(lm)}=R_{iklm}$, containing
$6^2=36$ potentially independent entries.

Pair exchange, $R_{iklm}=R_{lmik}$, says that $C$ is symmetric. Its six
diagonal entries and $6\cdot5/2=15$ entries above the diagonal determine
all the others. The number is now $6+15=21$.

Consider the cyclic identity. Its left side is unchanged by the cyclic
permutation $(k,l,m)\to(l,m,k)$, which merely permutes its three terms.
So if two of $k,l,m$ coincide, we may take them to be $k=l$. The left
side then has the form
\[
R_{ikkm}+R_{ikmk}+R_{imkk}=R_{ikkm}-R_{ikkm}+0=0.
\]
If $i$ coincides with one of $k,l,m$, the same cyclic freedom lets us
take $i=k$. Then
\[
R_{iilm}+R_{ilmi}+R_{imil}
=0-R_{ilim}+R_{ilim}=0,
\]
because $R_{ilmi}=-R_{ilim}$ by antisymmetry in the last pair, and
$R_{imil}=R_{ilim}$ by pair exchange.
Thus a new restriction is possible only when all four indices differ.
For the four available indices it is
\[
R_{0123}+R_{0231}+R_{0312}=0,
\quad\hbox{or}\quad
\boxed{C_{01,23}-C_{02,13}+C_{03,12}=0}.
\]
All other orderings give this same equation, possibly with the opposite
sign. To see why, set
$B_{iklm}=R_{iklm}+R_{ilmk}+R_{imkl}$. Interchanging $k,l$ changes all
three terms' signs, by antisymmetry in the last pair:
\[
B_{ilkm}=R_{ilkm}+R_{ikml}+R_{imlk}
=-R_{ilmk}-R_{iklm}-R_{imkl}=-B_{iklm}.
\]
Interchanging $l,m$ likewise reverses all three signs:
\[
B_{ikml}=R_{ikml}+R_{imlk}+R_{ilkm}
=-R_{iklm}-R_{imkl}-R_{ilmk}=-B_{iklm}.
\]
Interchanging $i,k$ gives
\[
B_{kilm}=R_{kilm}+R_{klmi}+R_{kmil}
=-R_{iklm}-R_{imkl}-R_{ilmk}=-B_{iklm}.
\]
Here $R_{kilm}=-R_{iklm}$, while pair exchange followed by antisymmetry
gives $R_{klmi}=R_{mikl}=-R_{imkl}$ and $R_{kmil}=R_{ilkm}=-R_{ilmk}$.
Every permutation of the four slots is a product of interchanges of
neighboring slots, so $B$ is completely antisymmetric. Its only
independent entry in four dimensions is $B_{0123}$. The displayed relation is
nontrivial: it links three different entries above the diagonal of the
symmetric array $C$, which are independent among the 21 parameters. It
therefore removes exactly one of them. The final number is
\[
\boxed{21-1=20}.
\]

\Needspace{5\baselineskip}
\par\bigskip
\noindent\textbf{4.1.6}\quad $T_{iklm}$ is antisymmetric with respect to all
pairs of indices. How many independent components has it (in 3-D space)?
\solution
Four indices chosen from the three values $1,2,3$ must include a
repeated value. Interchanging the two slots containing that repeated
value leaves a component unchanged. Complete antisymmetry says that the
same interchange changes its sign. Therefore each component obeys
$T_{iklm}=-T_{iklm}$ and is zero. The number of independent components is
\[
\boxed{0}.
\]
Equivalently, the space of completely antisymmetric rank-four tensors
in three dimensions has dimension $\binom34=0$.

\Needspace{5\baselineskip}
\par\bigskip
\noindent\textbf{4.1.7}\quad If $T_{\ldots i}$ is a tensor of rank $n$, show
that $\partial T_{\ldots i}/\partial x^{j}$ is a tensor of rank $n+1$
(Cartesian coordinates).

\emph{Note.}~In non-Cartesian coordinate systems the coefficients $a_{ij}$
are, in general, functions of the coordinates, and the derivatives \erratum{of} the
components of a tensor of rank $n$ do not form a tensor except in the special
case $n=0$. In this case the derivative does yield a covariant vector (tensor
of rank $1$).
\solution
A change between orthonormal Cartesian systems has the form
$x'^i=Q_{ia}x^a+s^i$, where $Q$ and the shift $s^i$ are constant and
$Q^{\mathsf T}Q=I$. Multiplying by $Q_{ib}$, summing over $i$, and using
$Q_{ib}Q_{ia}=\delta_{ba}$ gives the inverse map
$x^b=Q_{ib}(x'^i-s^i)$, so $\partial x^b/\partial x'^j=Q_{jb}$. The
chain rule then gives
\[
\partial'_j\equiv\frac{\partial}{\partial x'^j}
=\frac{\partial x^b}{\partial x'^j}\frac{\partial}{\partial x^b}
=Q_{jb}\partial_b.
\]
For a rank-$n$ tensor, its Cartesian component law is
\[
T'_{i_1\cdots i_n}=Q_{i_1a_1}\cdots Q_{i_na_n}T_{a_1\cdots a_n}.
\]
Differentiate this equation. The $Q$'s are constant, so the product rule
produces no derivatives of them:
\[
\partial'_jT'_{i_1\cdots i_n}
=Q_{i_1a_1}\cdots Q_{i_na_n}Q_{jb}
  \partial_bT_{a_1\cdots a_n}.
\]
There is one transformation factor for each of the original $n$ indices
and one more for $j$. This is precisely a tensor law of rank $n+1$.
For general coordinates the Jacobian factors depend on position, and
differentiating them gives extra terms. For a covariant vector, for
instance,
\[
\partial'_jV'_i
=\frac{\partial}{\partial x'^j}
 \Bigl(\frac{\partial x^a}{\partial x'^i}V_a\Bigr)
=\frac{\partial x^a}{\partial x'^i}\frac{\partial x^b}{\partial x'^j}
 \partial_bV_a
 +\frac{\partial^2x^a}{\partial x'^j\partial x'^i}V_a,
\]
and the last term spoils the tensor law. A scalar has no such factors,
so its ordinary first derivative remains a covector in general coordinates.
\erratum{The word ``of'' in the Note is supplied: the printed sentence
lacks it (``the derivatives the components''), but it is the
derivatives \emph{of} the components that fail to form a tensor.}

\Needspace{5\baselineskip}
\par\bigskip
\noindent\textbf{4.1.8}\quad If $T_{ijk\ldots}$ is a tensor of rank $n$, show
that $\sum_{j}\partial T_{ijk\ldots}/\partial x^{j}$ is a tensor of rank
$n-1$ (Cartesian coordinates).
\solution
Let $D_{ik\cdots}=\partial_jT_{ijk\cdots}$, with the repeated
$j$ summed. Use the same constant orthogonal matrix as in Exercise 4.1.7:
\begin{align*}
D'_{ik\cdots}
&=\partial'_jT'_{ijk\cdots}\\
&=Q_{jd}\partial_d
  \bigl(Q_{ia}Q_{jb}Q_{kc}\cdots T_{abc\cdots}\bigr)\\
&=Q_{ia}Q_{kc}\cdots
  \left(\sum_jQ_{jd}Q_{jb}\right)\partial_dT_{abc\cdots}\\
&=Q_{ia}Q_{kc}\cdots\delta_{db}\partial_dT_{abc\cdots}\\
&=Q_{ia}Q_{kc}\cdots\partial_bT_{abc\cdots}
 =Q_{ia}Q_{kc}\cdots D_{ac\cdots}.
\end{align*}
The constant $Q$'s were taken outside the derivative, and the fourth line
uses orthogonality,
$\sum_jQ_{jd}Q_{jb}=(Q^{\mathsf T}Q)_{db}=\delta_{db}$.
Differentiation introduces one index and contraction removes two, leaving
$n+1-2=n-1$ free indices. The displayed law verifies that the result is a
rank-$(n-1)$ tensor. In the rank-one case the same calculation gives the
scalar divergence $\partial_jT_j$.

\Needspace{5\baselineskip}
\par\bigskip
\noindent\textbf{4.1.9}\quad The operator
\[
\nabla^{2}-\frac{1}{c^{2}}\frac{\partial^{2}}{\partial t^{2}}
\]
may be written as
\[
\sum_{i=1}^{4}\frac{\partial^{2}}{\partial x_{i}^{2}},
\]
using $x_{4}=ict$. This is the 4-D Laplacian, sometimes called the
d'Alembertian and denoted by $\square^{2}$. Show that it is a scalar operator,
that is, invariant under Lorentz transformations, represented here by
\erratum{complex orthogonal transformations} in the space of vectors $(x_{1},x_{2},x_{3},x_{4})$.
\solution
Since $x_4=ict$, differentiation at fixed spatial coordinates gives
\[
\frac{\partial}{\partial x_4}=\frac{1}{ic}\frac{\partial}{\partial t},
\qquad
\frac{\partial^2}{\partial x_4^2}
=-\frac{1}{c^2}\frac{\partial^2}{\partial t^2}.
\]
Thus $\sum_{a=1}^4\partial_a^2$ is the stated operator. In these complex
coordinates a Lorentz transformation $x'=Qx$ is represented by a
constant complex orthogonal matrix $Q$: it preserves
$x_1^2+x_2^2+x_3^2+x_4^2=\mathbf{x}^2-c^2t^2$, that is,
$x^{\mathsf T}Mx=0$ for every event $x$, where $M=Q^{\mathsf T}Q-I$ is
symmetric. Sums of events are events, so polarization, as in
Exercise 4.1.10, gives
\[
u^{\mathsf T}Mv=\tfrac12\{(u+v)^{\mathsf T}M(u+v)-u^{\mathsf T}Mu
-v^{\mathsf T}Mv\}=0
\]
for any two events $u,v$. Events have real $x_1,x_2,x_3$ and imaginary
$x_4$, so $u,v$ may be taken among the unit columns $e_1,e_2,e_3$ and
$ie_4$; this gives $M_{bd}=0$, $iM_{b4}=0$ ($b,d\le3$) and $-M_{44}=0$.
Hence $M=0$, that is, $Q^{\mathsf T}Q=I$. For example, the boost
$x_1'=\gamma(x_1-\beta ct)$, $t'=\gamma(t-\beta x_1/c)$, with
$\gamma=(1-\beta^2)^{-1/2}$, reads
\[
x_1'=\gamma(x_1+i\beta x_4),\quad x_4'=\gamma(x_4-i\beta x_1),
\qquad
\begin{pmatrix}Q_{11}&Q_{14}\\Q_{41}&Q_{44}\end{pmatrix}
=\begin{pmatrix}\gamma&i\beta\gamma\\-i\beta\gamma&\gamma\end{pmatrix},
\]
and this block satisfies $Q^{\mathsf T}Q=I$: the diagonal entries of
the product are $\gamma^2+(i\beta\gamma)^2=\gamma^2(1-\beta^2)=1$, and
the off-diagonal ones are $i\beta\gamma^2-i\beta\gamma^2=0$.
These are complex rotations, not arbitrary real
Euclidean rotations mixing the physical time and space coordinates.
Since $x=Q^{\mathsf T}x'$, that is, $x_b=Q_{ab}x'_a$, the chain rule gives
$\partial'_a=(\partial x_b/\partial x'_a)\partial_b=Q_{ab}\partial_b$.
Hence, acting on a scalar function,
\begin{align*}
\sum_a\partial'^2_a
&=\sum_{a,b,d}Q_{ab}Q_{ad}\partial_b\partial_d\\
&=\sum_{b,d}(Q^{\mathsf T}Q)_{bd}\partial_b\partial_d
 =\sum_{b,d}\delta_{bd}\partial_b\partial_d
 =\sum_b\partial_b^2.
\end{align*}
This proves the operator's Lorentz invariance. No derivative of $Q$
appears because inertial Lorentz transformations are linear.

\Needspace{5\baselineskip}
\par\bigskip
\noindent\textbf{4.1.10}\quad The double summation $K_{ij}A_{i}B_{j}$ is
invariant for any two vectors $A_{i}$ and $B_{j}$. Prove that $K_{ij}$ is a
second-rank tensor.

\emph{Note.}~In the form $ds^{2}(\text{invariant})=g_{ij}\,dx^{i}\,dx^{j}$,
this result shows that the \erratum{symmetric} matrix $g_{ij}$ is a tensor.
\solution
Use orthonormal Cartesian coordinates, so vectors transform as
$A'_p=Q_{pi}A_i$ and $B'_q=Q_{qj}B_j$. Invariance of the given scalar
means
\[
K_{ij}A_iB_j
=K'_{pq}A'_pB'_q
=K'_{pq}Q_{pi}Q_{qj}A_iB_j.
\]
Choose $A$ and $B$ successively to be each pair of coordinate unit
vectors, $A_i=\delta_{im}$ and $B_j=\delta_{jn}$; every such choice is
permitted. The sums over $i$ and $j$ then collapse, leaving
$K_{mn}=Q_{pm}Q_{qn}K'_{pq}$ for all $m,n$. Renaming $m,n$ as $i,j$,
\[
K_{ij}=Q_{pi}Q_{qj}K'_{pq}.
\]
Multiply by $Q_{ri}Q_{sj}$ and sum $i,j$. Orthogonality gives
\[
Q_{ri}Q_{sj}K_{ij}
=(Q_{ri}Q_{pi})(Q_{sj}Q_{qj})K'_{pq}
=\delta_{rp}\delta_{sq}K'_{pq}=K'_{rs}.
\]
Thus $K'_{rs}=Q_{ri}Q_{sj}K_{ij}$, the second-rank tensor law.

\erratum{The line-element application uses the symmetry of the metric.
A quadratic expression alone cannot determine an antisymmetric part of
$g_{ij}$, because that part contracts to zero with $dx^i dx^j$.}
For a symmetric metric, polarization recovers the bilinear expression
from the invariant quadratic form $q(U)=g_{ij}U^iU^j$:
\[
g_{ij}U^iV^j=\frac12\{q(U+V)-q(U)-q(V)\}.
\]
To see the polarization step explicitly, expand
\[
q(U+V)=q(U)+q(V)+g_{ij}U^iV^j+g_{ij}V^iU^j
=q(U)+q(V)+2g_{ij}U^iV^j,
\]
where symmetry and relabeling $i,j$ equate the cross terms.
The left side of the polarization formula is consequently invariant for
independent $U,V$. For a general nonsingular coordinate change, write
$U^i=J^i{}_aU'^a$, $V^j=J^j{}_bV'^b$, where
$J^i{}_a=\partial x^i/\partial x'^a$. Then
\[
g'_{ab}U'^aV'^b
=g_{ij}J^i{}_aJ^j{}_bU'^aV'^b.
\]
Taking $U',V'$ to be each pair of coordinate unit vectors yields
$g'_{ab}=J^i{}_aJ^j{}_b g_{ij}$, the covariant metric transformation law.

\Needspace{5\baselineskip}
\par\bigskip
\noindent\textbf{4.1.11}\quad The equation $K_{ij}A_{jk}=B_{ik}$ holds for all
orientations of the coordinate system. If $A$ and $B$ are arbitrary
second-rank tensors, show that $K$ is a second-rank tensor also.
\solution
Write the contraction as the matrix equation $KA=B$. Under a
Cartesian rotation $Q$, the two known tensors transform as
$A'=QAQ^{\mathsf T}$ and $B'=QBQ^{\mathsf T}$. Substituting $B=KA$ and
inserting $Q^{\mathsf T}Q=I$ between $K$ and $A$,
\[
B'=QKAQ^{\mathsf T}
  =QK(Q^{\mathsf T}Q)AQ^{\mathsf T}
  =(QKQ^{\mathsf T})(QAQ^{\mathsf T})
  =(QKQ^{\mathsf T})A'.
\]
But the hypothesis states that the equation has the same form in the
new coordinates: $K'A'=B'$. Subtracting gives
\[
(K'-QKQ^{\mathsf T})A'=0.
\]
This must hold for every second-rank $A'$. In particular, set $A'=I$;
then $K'-QKQ^{\mathsf T}=0$. Hence
\[
\boxed{K'=QKQ^{\mathsf T}},
\]
which is the required tensor law. The word ``arbitrary'' is essential:
if the equation were supplied only for one singular $A$, it would not
in general determine $K$; for example $A=B=0$ imposes no condition.

\Needspace{9\baselineskip}
\subsection*{4.2 Pseudotensors, Dual Tensors}

\Needspace{5\baselineskip}
\par\bigskip
\noindent\textbf{4.2.1}\quad An antisymmetric square array is given by
\[
\begin{pmatrix}0&C_{3}&-C_{2}\\-C_{3}&0&C_{1}\\C_{2}&-C_{1}&0\end{pmatrix}
=\begin{pmatrix}0&C_{12}&C_{13}\\-C_{12}&0&C_{23}\\
-C_{13}&-C_{23}&0\end{pmatrix},
\]
where $(C_{1},C_{2},C_{3})$ form a pseudovector. Assuming that the relation
\[
C_{i}=\frac{1}{2!}\epsilon_{ijk}C_{jk}
\]
holds in all \erratum{orthonormal Cartesian coordinate systems}, prove that $C_{jk}$ is a tensor. (This is
another form of the quotient theorem.)
\solution
\erratum{The numerical Levi-Civita symbol in this question uses
orthonormal Cartesian axes. General curvilinear coordinates require the
metric-dependent volume tensor, rather than the same numerical symbol.}
First invert the stated relation: multiply it by $\epsilon_{\ell mi}$
and sum over $i$. The contraction identity
\[
\epsilon_{\ell mi}\epsilon_{ijk}
=\delta_{\ell j}\delta_{mk}-\delta_{\ell k}\delta_{mj}
\]
gives
\[
\epsilon_{\ell mi}C_i
=\frac12(\delta_{\ell j}\delta_{mk}-\delta_{\ell k}\delta_{mj})C_{jk}
=\frac12(C_{\ell m}-C_{m\ell})=C_{\ell m}.
\]
Here antisymmetry was essential in the final step. Thus
$C_{jk}=\epsilon_{jki}C_i$, in every orthonormal Cartesian system.

For an orthogonal change of Cartesian axes $Q$, put $d=\det Q=\pm1$.
A pseudovector obeys $C'_i=dQ_{ia}C_a$. For the numerical Levi-Civita
symbols,
$Q_{jb}Q_{kc}Q_{i\ell}\epsilon_{bc\ell}=(\det Q)\epsilon_{jki}$.
Indeed, by the expansion $\det A=\epsilon_{bc\ell}A_{1b}A_{2c}A_{3\ell}$,
the left side is the determinant of the matrix whose rows are rows
$j,k,i$ of $Q$: for $(j,k,i)=(1,2,3)$ it is $\det Q$, an interchange of
two rows reverses its sign, and a repeated row gives zero, exactly as
for $(\det Q)\epsilon_{jki}$. Multiplying by $d$ and using $d^2=1$ yields
\[
\epsilon_{jki}=dQ_{jb}Q_{kc}Q_{i\ell}\epsilon_{bc\ell}.
\]
Substituting both identities into the inverted relation gives
\begin{align*}
C'_{jk}
&=\epsilon_{jki}C'_i\\
&=d^2Q_{jb}Q_{kc}Q_{i\ell}Q_{ia}\epsilon_{bc\ell}C_a\\
&=Q_{jb}Q_{kc}\delta_{\ell a}\epsilon_{bc\ell}C_a\\
&=Q_{jb}Q_{kc}C_{bc}.
\end{align*}
We used $d^2=1$, $\sum_iQ_{i\ell}Q_{ia}=\delta_{\ell a}$, and, in the
last line, the inverted relation $\epsilon_{bca}C_a=C_{bc}$ in the
original axes.
This is the ordinary second-rank tensor law, including reflections.

\Needspace{5\baselineskip}
\par\bigskip
\noindent\textbf{4.2.2}\quad Show that the vector product is unique to 3-D
space, that is, only in three dimensions can we establish a
\erratum{linear} one-to-one correspondence between the components of an
antisymmetric tensor (second-rank) and the components of a vector.
\solution
In $n$ dimensions, an antisymmetric tensor has zero diagonal
entries, and each entry below the diagonal is minus the corresponding
entry above it. Its independent components are therefore indexed by
unordered pairs of distinct indices, whose number is
\[
\binom n2=\frac{n(n-1)}2.
\]
A vector has $n$ independent components. A linear one-to-one
correspondence between the full component spaces requires equal
dimensions:
\[
\frac{n(n-1)}2=n
\quad\Longrightarrow\quad n(n-3)=0.
\]
For a positive space dimension this forces $n=3$. In three dimensions
the correspondence actually exists: $C_i=\tfrac12\epsilon_{ijk}C_{jk}$
and $C_{jk}=\epsilon_{jki}C_i$ are inverse maps. Exercise 4.2.1 showed
that the second undoes the first. In the other order, start from $C_j$
and form $C_{mi}=\epsilon_{mij}C_j=\epsilon_{ijm}C_j$. Setting $k=m$ in
the contraction identity of Exercise 4.2.1 and summing gives
$\epsilon_{\ell mi}\epsilon_{ijm}=3\delta_{\ell j}-\delta_{\ell j}
=2\delta_{\ell j}$, since $\delta_{mm}=3$; hence
$\tfrac12\epsilon_{\ell mi}C_{mi}
=\tfrac12\epsilon_{\ell mi}\epsilon_{ijm}C_j=C_\ell$.
The associated vector is a pseudovector when orientation reversals are
included.

\erratum{The word ``linear'' is added to the statement because, as mere
sets, $\mathbb R^{n(n-1)/2}$ and $\mathbb R^n$ can be put in one-to-one
correspondence for every $n\ge2$. The uniqueness of the vector product
is meant in the sense of this linear correspondence with all
antisymmetric second-rank tensors; the dimension count does not exclude
other, differently defined products bearing the name cross product.}

\Needspace{5\baselineskip}
\par\bigskip
\noindent\textbf{4.2.3}\quad Write $\nabla\cdot\nabla\times\mathbf{A}$ and
$\nabla\times\nabla\varphi$ in tensor (index) notation in $\mathbb{R}^{3}$ so
that it becomes obvious that each expression vanishes.
\erratum{Assume that the fields have continuous second derivatives.}
\par\nobreak\hfill\emph{ANS.}\enspace
$\begin{aligned}
\nabla\cdot\nabla\times\mathbf{A}
&=\epsilon_{ijk}\frac{\partial}{\partial x^{i}}\frac{\partial}{\partial x^{j}}
A^{k}\\
(\nabla\times\nabla\varphi)_{i}
&=\epsilon_{ijk}\frac{\partial}{\partial x^{j}}\frac{\partial}{\partial x^{k}}
\varphi.
\end{aligned}$
\solution
Assume the fields have continuous second derivatives so that mixed
partial derivatives commute. With repeated indices summed,
\[
\nabla\cdot(\nabla\times\mathbf A)
=\partial_i(\epsilon_{ijk}\partial_jA^k)
=\epsilon_{ijk}\partial_i\partial_jA^k.
\]
Interchanging the dummy indices $i,j$ changes $\epsilon_{ijk}$ to its
negative but leaves $\partial_i\partial_jA^k$ unchanged. More explicitly,
relabeling gives
$\epsilon_{ijk}\partial_i\partial_jA^k=\epsilon_{jik}\partial_j\partial_iA^k
=-\epsilon_{ijk}\partial_j\partial_iA^k$, and the average of these two
forms is
\[
\epsilon_{ijk}\partial_i\partial_jA^k
=\frac12\epsilon_{ijk}
 (\partial_i\partial_j-\partial_j\partial_i)A^k=0.
\]
Likewise,
\[
(\nabla\times\nabla\varphi)_i
=\epsilon_{ijk}\partial_j\partial_k\varphi
=\frac12\epsilon_{ijk}
 (\partial_j\partial_k-\partial_k\partial_j)\varphi=0.
\]
Both identities are contractions of an antisymmetric pair of indices
with a symmetric pair.

\Needspace{5\baselineskip}
\par\bigskip
\noindent\textbf{4.2.4}\quad Verify that each of the following fourth-rank
tensors is isotropic, that is, that it has the same form independent of any
rotation of the coordinate systems.
\begin{enumerate}
\item[(a)] $A_{jl}^{ik}=\delta_{j}^{i}\delta_{l}^{k}$,
\item[(b)] $B_{kl}^{ij}=\delta_{k}^{i}\delta_{l}^{j}
  +\delta_{l}^{i}\delta_{k}^{j}$,
\item[(c)] $C_{kl}^{ij}=\delta_{k}^{i}\delta_{l}^{j}
  -\delta_{l}^{i}\delta_{k}^{j}$.
\end{enumerate}
\solution
The mixed identity tensor is invariant. Indeed, under a rotation,
\[
\delta'{}^i_j
=Q_{ia}Q_{jb}\delta^a_b
=\sum_aQ_{ia}Q_{ja}
=(QQ^{\mathsf T})_{ij}=\delta^i_j.
\]
For (a), apply the four-index transformation law and group the factors
associated with each Kronecker delta:
\begin{align*}
A'{}^{ik}_{jl}
&=Q_{ia}Q_{kc}Q_{jb}Q_{ld}\delta^a_b\delta^c_d\\
&=\left(\sum_aQ_{ia}Q_{ja}\right)
  \left(\sum_cQ_{kc}Q_{lc}\right)
=\delta^i_j\delta^k_l=A^{ik}_{jl}.
\end{align*}
The same calculation with the indices arranged as in (b) gives
\[
(\delta^i_k\delta^j_l)'=\delta^i_k\delta^j_l,
\qquad
(\delta^i_l\delta^j_k)'=\delta^i_l\delta^j_k.
\]
The second pairs the indices crosswise; with upper indices $i,j$ and
lower indices $k,l$ its transformation reads
\[
Q_{ia}Q_{jc}Q_{kb}Q_{ld}\delta^a_d\delta^c_b
=\Bigl(\sum_aQ_{ia}Q_{la}\Bigr)\Bigl(\sum_cQ_{jc}Q_{kc}\Bigr)
=\delta^i_l\delta^j_k.
\]
Since transformation is linear, adding the two invariance equations proves
$B'{}^{ij}_{kl}=B^{ij}_{kl}$, and subtracting them proves
$C'{}^{ij}_{kl}=C^{ij}_{kl}$. Thus all three tensors are isotropic.
The argument in fact also allows orthogonal reflections.

\Needspace{5\baselineskip}
\par\bigskip
\noindent\textbf{4.2.5}\quad Show that the two-index Levi-Civita symbol
$\epsilon_{ij}$ is a second-rank pseudotensor (in two-dimensional [2-D]
space). Does this contradict the uniqueness of $\delta_{j}^{i}$
(Exercise 4.1.4)?
\solution
Use the orientation convention $\epsilon_{12}=1$, so
\[
E=(\epsilon_{ij})=\begin{pmatrix}0&1\\-1&0\end{pmatrix}.
\]
For $Q=\begin{pmatrix}a&b\\c&d\end{pmatrix}$, direct multiplication gives
\[
QEQ^{\mathsf T}
=\begin{pmatrix}-b&a\\-d&c\end{pmatrix}
 \begin{pmatrix}a&c\\b&d\end{pmatrix}
=\begin{pmatrix}0&ad-bc\\bc-ad&0\end{pmatrix}
=(\det Q)E.
\]
In particular, for orthogonal $Q$ the pseudotensor law is
\[
E'=(\det Q)QEQ^{\mathsf T}=(\det Q)^2E=E.
\]
Thus the same numerical symbol in every oriented coordinate convention
is a second-rank pseudotensor. A proper rotation has $\det Q=1$; a
reflection has $\det Q=-1$ and contributes the extra sign appropriate
to a pseudotensor.

There is no contradiction with Exercise 4.1.4, which concerns
\emph{three} dimensions. In two dimensions $E$ is a nonzero
antisymmetric matrix invariant under all proper rotations; the rotations
about three distinct axes used in that exercise are unavailable.
Moreover, $E$ is not invariant under the ordinary tensor transformation
law for reflections, whereas the identity is.

\Needspace{5\baselineskip}
\par\bigskip
\noindent\textbf{4.2.6}\quad Represent $\epsilon_{ij}$ by a $2\times2$ matrix,
and using the $2\times2$ rotation matrix of Eq.~(3.23), show that
$\epsilon_{ij}$ is invariant under \erratum{proper orthogonal} similarity transformations.
\solution
\erratum{Here ``orthogonal similarity transformations'' must be
restricted to proper rotations ($\det Q=1$). For general orthogonal $Q$
the correct relation is $QEQ^{-1}=(\det Q)E$.}

With $c=\cos\phi$, $s=\sin\phi$, the rotation matrix in the question is
\[
R=\begin{pmatrix}c&s\\-s&c\end{pmatrix},\qquad
R^{-1}=R^{\mathsf T}=\begin{pmatrix}c&-s\\s&c\end{pmatrix}.
\]
Since $E=\begin{pmatrix}0&1\\-1&0\end{pmatrix}$, multiply explicitly:
\[
RER^{-1}
=\begin{pmatrix}-s&c\\-c&-s\end{pmatrix}
 \begin{pmatrix}c&-s\\s&c\end{pmatrix}
=\begin{pmatrix}0&s^2+c^2\\-(c^2+s^2)&0\end{pmatrix}
=E.
\]
This proves the desired rotational invariance. To check the necessary
qualification, take the reflection $Q=\operatorname{diag}(1,-1)$.
Then
\[
QEQ^{-1}=\begin{pmatrix}0&-1\\1&0\end{pmatrix}=-E.
\]
Under the \emph{pseudotensor} law its additional factor $\det Q=-1$
restores $E$, exactly as in Exercise 4.2.5.

\Needspace{5\baselineskip}
\par\bigskip
\noindent\textbf{4.2.7}\quad Given
$A_{k}=\frac{1}{2}\epsilon_{ijk}B^{ij}$ with $B^{ij}=-B^{ji}$, antisymmetric,
show that
\[
B^{mn}=\epsilon^{mnk}A_{k}.
\]
\solution
Multiply the given equation by $\epsilon^{mnk}$ and sum $k$:
\begin{align*}
\epsilon^{mnk}A_k
&=\frac12\epsilon^{mnk}\epsilon_{ijk}B^{ij}\\
&=\frac12(\delta^m_i\delta^n_j-\delta^m_j\delta^n_i)B^{ij}\\
&=\frac12(B^{mn}-B^{nm})\\
&=\frac12(B^{mn}+B^{mn})=B^{mn}.
\end{align*}
The second line is the contraction identity of Exercise 4.2.1, applied
to $\epsilon^{mnk}\epsilon_{kij}$ after the cyclic rearrangement
$\epsilon_{ijk}=\epsilon_{kij}$ (in Cartesian coordinates upper and
lower indices are equivalent), and the last line uses $B^{nm}=-B^{mn}$.
This is the stated inverse relation. It also checks the factor $1/2$:
without it the two antisymmetric terms would produce $2B^{mn}$.

\Needspace{9\baselineskip}
\subsection*{4.3 Tensors in General Coordinates}

\Needspace{5\baselineskip}
\par\bigskip
\noindent\textbf{4.3.1}\quad For the special case of 3-D space
($\boldsymbol{\varepsilon}_{1},\boldsymbol{\varepsilon}_{2},
\boldsymbol{\varepsilon}_{3}$ defining a right-handed coordinate system, not
necessarily orthogonal), show that
\[
\boldsymbol{\varepsilon}^{i}
=\frac{\boldsymbol{\varepsilon}_{j}\times\boldsymbol{\varepsilon}_{k}}
{\boldsymbol{\varepsilon}_{j}\times\boldsymbol{\varepsilon}_{k}\cdot
\boldsymbol{\varepsilon}_{i}},\quad i,j,k=1,2,3\text{ and cyclic
permutations.}
\]
\emph{Note.}~These contravariant basis vectors
$\boldsymbol{\varepsilon}^{i}$ define the reciprocal lattice space of
Example 3.2.1.
\solution
The reciprocal basis is characterized by
$\boldsymbol\varepsilon^i\cdot\boldsymbol\varepsilon_\ell
=\delta^i_\ell$: with $\boldsymbol\varepsilon^i=\nabla q^i$ and
$\boldsymbol\varepsilon_\ell=\partial\mathbf r/\partial q^\ell$, the chain
rule gives
$\nabla q^i\cdot\partial\mathbf r/\partial q^\ell
=\partial q^i/\partial q^\ell=\delta^i_\ell$. Let
\[
D=(\boldsymbol\varepsilon_1\times\boldsymbol\varepsilon_2)
 \cdot\boldsymbol\varepsilon_3.
\]
Linear independence implies $D\ne0$, and the specified right-handedness
makes $D>0$. Cyclic permutations preserve this scalar triple product,
so for cyclic $(i,j,k)$ the denominator in the formula is
$(\boldsymbol\varepsilon_j\times\boldsymbol\varepsilon_k)
\cdot\boldsymbol\varepsilon_i=D$, and the right side of the formula is
\[
\mathbf b^i=\frac{\boldsymbol\varepsilon_j\times
                         \boldsymbol\varepsilon_k}{D}.
\]
The cross product is perpendicular to both its factors, so
\[
\mathbf b^i\cdot\boldsymbol\varepsilon_j=0,\qquad
\mathbf b^i\cdot\boldsymbol\varepsilon_k=0.
\]
Its remaining scalar product is
\[
\mathbf b^i\cdot\boldsymbol\varepsilon_i
=\frac{(\boldsymbol\varepsilon_j\times\boldsymbol\varepsilon_k)
            \cdot\boldsymbol\varepsilon_i}{D}=1.
\]
These are exactly the three reciprocal-basis conditions. They determine
a vector uniquely: the difference of any two solutions would be
orthogonal to three linearly independent vectors and hence zero.
Therefore $\mathbf b^i=\boldsymbol\varepsilon^i$, proving the formula.

\Needspace{5\baselineskip}
\par\bigskip
\noindent\textbf{4.3.2}\quad If the covariant vectors
$\boldsymbol{\varepsilon}_{i}$ are orthogonal, show that
\begin{enumerate}
\item[(a)] $g_{ij}$ is diagonal,
\item[(b)] $g^{ii}=1/g_{ii}$ (no summation),
\item[(c)] $|\boldsymbol{\varepsilon}^{i}|=1/|\boldsymbol{\varepsilon}_{i}|$.
\end{enumerate}
\solution
Put $h_i=|\boldsymbol\varepsilon_i|>0$. Orthogonality means
\[
g_{ij}=\boldsymbol\varepsilon_i\cdot\boldsymbol\varepsilon_j
=\begin{cases}h_i^2,&i=j,\\0,&i\ne j.\end{cases}
\]
Thus (a) the covariant metric is diagonal. The contravariant metric
$g^{ij}=\boldsymbol\varepsilon^i\cdot\boldsymbol\varepsilon^j$ is its
inverse, $\sum_jg^{ij}g_{jk}=\delta^i_k$ (Exercise 4.3.3). For a diagonal
$g_{jk}$ only the term $j=k$ survives, so this reads
$g^{ik}h_k^2=\delta^i_k$ (no sum), and
\[
g^{ik}=0\quad(i\ne k),\qquad
\boxed{g^{ii}=h_i^{-2}=1/g_{ii}}
\quad\text{(no sum)}.
\]
This proves (b). To obtain (c), observe that
$\boldsymbol\varepsilon_i/h_i^2$ has scalar product $1$ with
$\boldsymbol\varepsilon_i$ and $0$ with all the other basis vectors.
By the uniqueness argument of Exercise 4.3.1, it is consequently the
reciprocal vector:
\[
\boldsymbol\varepsilon^i
=\frac{\boldsymbol\varepsilon_i}{h_i^2},\qquad
\boxed{|\boldsymbol\varepsilon^i|
=\frac{h_i}{h_i^2}=\frac1{|\boldsymbol\varepsilon_i|}}.
\]

\Needspace{5\baselineskip}
\par\bigskip
\noindent\textbf{4.3.3}\quad Prove that
$(\boldsymbol{\varepsilon}^{i}\cdot\boldsymbol{\varepsilon}^{j})
(\boldsymbol{\varepsilon}_{j}\cdot\boldsymbol{\varepsilon}_{k})
=\delta_{k}^{i}$.
\solution
The dual-basis expansion of any vector $\mathbf v$ is
\[
\mathbf v=\sum_j(\boldsymbol\varepsilon_j\cdot\mathbf v)
                  \boldsymbol\varepsilon^j.
\]
To verify it, take the scalar product with each
$\boldsymbol\varepsilon_k$; the right side gives
$\sum_j(\boldsymbol\varepsilon_j\cdot\mathbf v)\delta^j_k
=\boldsymbol\varepsilon_k\cdot\mathbf v$. A vector is determined by
these three scalar products. Apply the expansion to
$\mathbf v=\boldsymbol\varepsilon_k$ and then dot with
$\boldsymbol\varepsilon^i$:
\[
\boldsymbol\varepsilon^i\cdot\boldsymbol\varepsilon_k
=\sum_j(\boldsymbol\varepsilon_j\cdot\boldsymbol\varepsilon_k)
       (\boldsymbol\varepsilon^i\cdot\boldsymbol\varepsilon^j)
=\delta^i_k.
\]
This is the desired identity, also written $g^{ij}g_{jk}=\delta^i_k$.

\Needspace{5\baselineskip}
\par\bigskip
\noindent\textbf{4.3.4}\quad Show that
$\Gamma_{jk}^{m}=\Gamma_{kj}^{m}$.
\solution
For a smooth coordinate map $\mathbf r(q^1,q^2,q^3)$,
\[
\boldsymbol\varepsilon_j=\frac{\partial\mathbf r}{\partial q^j},
\qquad
\frac{\partial\boldsymbol\varepsilon_j}{\partial q^k}
=\Gamma^m_{jk}\boldsymbol\varepsilon_m.
\]
Commuting the mixed second derivatives of $\mathbf r$ gives
\[
\Gamma^m_{jk}\boldsymbol\varepsilon_m
=\frac{\partial^2\mathbf r}{\partial q^k\partial q^j}
=\frac{\partial^2\mathbf r}{\partial q^j\partial q^k}
=\Gamma^m_{kj}\boldsymbol\varepsilon_m.
\]
Linear independence of the basis forces equality of each coefficient:
$\boxed{\Gamma^m_{jk}=\Gamma^m_{kj}}$. Equivalently, take the scalar
product of both sides with $\boldsymbol\varepsilon^n$ and use
$\boldsymbol\varepsilon_m\cdot\boldsymbol\varepsilon^n=\delta^n_m$ to
obtain $\Gamma^n_{jk}=\Gamma^n_{kj}$. The symmetry rests on the basis
vectors $\boldsymbol\varepsilon_j=\partial\mathbf r/\partial q^j$ being
derivatives of a single position vector $\mathbf r$. Coefficients
defined in the same way for a basis not of this form need not be
symmetric: for the cylindrical unit vectors,
$\partial_\phi\hat{\mathbf e}_\rho=\hat{\mathbf e}_\phi$ but
$\partial_\rho\hat{\mathbf e}_\phi=0$, so the analogue of
$\Gamma^\phi_{\rho\phi}$ is $1$ while that of $\Gamma^\phi_{\phi\rho}$
is $0$.

\Needspace{5\baselineskip}
\par\bigskip
\noindent\textbf{4.3.5}\quad Derive the covariant and contravariant metric
tensors for circular cylindrical coordinates.
\solution
Choose $q^1=\rho$, $q^2=\phi$, $q^3=z$, with
\[
\mathbf r=(\rho\cos\phi,\rho\sin\phi,z).
\]
Differentiating gives the coordinate basis
\[
\boldsymbol\varepsilon_\rho=(\cos\phi,\sin\phi,0),\quad
\boldsymbol\varepsilon_\phi=(-\rho\sin\phi,\rho\cos\phi,0),\quad
\boldsymbol\varepsilon_z=(0,0,1).
\]
Its squared lengths are $1,\rho^2,1$. The only potentially nonzero
off-diagonal scalar product is
$\boldsymbol\varepsilon_\rho\cdot\boldsymbol\varepsilon_\phi
=-\rho\cos\phi\sin\phi+\rho\sin\phi\cos\phi=0$;
the other off-diagonal scalar products also vanish, because
$\boldsymbol\varepsilon_z$ has no $x,y$ components while
$\boldsymbol\varepsilon_\rho$ and $\boldsymbol\varepsilon_\phi$ have no
$z$ component. The basis is orthogonal, so by Exercise 4.3.2(b) the
contravariant metric has the reciprocal diagonal entries. Therefore
\[
\boxed{(g_{ij})=\begin{pmatrix}1&0&0\\0&\rho^2&0\\0&0&1\end{pmatrix},
\qquad
(g^{ij})=\begin{pmatrix}1&0&0\\0&\rho^{-2}&0\\0&0&1\end{pmatrix}}.
\]
As checks, their product is the identity and
$ds^2=g_{ij}\,dq^i dq^j=d\rho^2+\rho^2d\phi^2+dz^2$.
These coordinate formulas apply for $\rho>0$; at the axis the angular
coordinate is singular.

\Needspace{5\baselineskip}
\par\bigskip
\noindent\textbf{4.3.6}\quad Show that the covariant derivative of a covariant
vector is given by
\[
V_{i;j}\equiv\frac{\partial V_{i}}{\partial q^{j}}-V_{k}\Gamma_{ij}^{k}.
\]
\emph{Hint.}~Differentiate
\[
\boldsymbol{\varepsilon}^{i}\cdot\boldsymbol{\varepsilon}_{j}
=\delta_{j}^{i}.
\]
\solution
Differentiate the reciprocal-basis identity of the hint, written as
$\boldsymbol\varepsilon^i\cdot\boldsymbol\varepsilon_k=\delta^i_k$, with
respect to $q^j$, using
$\partial_j\boldsymbol\varepsilon_k=\Gamma^m_{kj}\boldsymbol\varepsilon_m$:
\[
0=\partial_j(\boldsymbol\varepsilon^i\cdot\boldsymbol\varepsilon_k)
=(\partial_j\boldsymbol\varepsilon^i)\cdot\boldsymbol\varepsilon_k
 +\boldsymbol\varepsilon^i\cdot
    (\Gamma^m_{kj}\boldsymbol\varepsilon_m).
\]
The last term is $\Gamma^i_{kj}$, so
$(\partial_j\boldsymbol\varepsilon^i)\cdot\boldsymbol\varepsilon_k
=-\Gamma^i_{kj}$. Expanding this derivative in the reciprocal basis,
$\mathbf w=(\mathbf w\cdot\boldsymbol\varepsilon_k)\boldsymbol\varepsilon^k$
(Exercise 4.3.3), therefore gives
\[
\partial_j\boldsymbol\varepsilon^i
=-\Gamma^i_{kj}\boldsymbol\varepsilon^k.
\]
Now express $\mathbf V$ using covariant components,
$\mathbf V=V_i\boldsymbol\varepsilon^i$. Differentiation gives
\begin{align*}
\partial_j\mathbf V
&=(\partial_jV_i)\boldsymbol\varepsilon^i
  +V_i\partial_j\boldsymbol\varepsilon^i\\
&=(\partial_jV_i)\boldsymbol\varepsilon^i
  -V_i\Gamma^i_{kj}\boldsymbol\varepsilon^k\\
&=(\partial_jV_i-V_k\Gamma^k_{ij})\boldsymbol\varepsilon^i.
\end{align*}
In the last line the dummy indices in the second term were relabeled,
$i\leftrightarrow k$.
The coefficients of the reciprocal basis are the covariant components
of the derivative, hence
\[
\boxed{V_{i;j}=\partial_jV_i-V_k\Gamma^k_{ij}}.
\]

\Needspace{5\baselineskip}
\par\bigskip
\noindent\textbf{4.3.7}\quad Verify that $V_{i;j}=g_{ik}V^{k}_{;j}$ by showing
that
\[
\frac{\partial V_{i}}{\partial q^{j}}-V_{k}\Gamma_{ij}^{k}
=g_{ik}\left[\frac{\partial V^{k}}{\partial q^{j}}
+V^{m}\Gamma_{mj}^{k}\right].
\]
\solution
Start from $V_i=g_{ik}V^k$. We first need the derivative of the
metric. Differentiate its basis-vector definition with the product rule
and use
$\partial_j\boldsymbol\varepsilon_i=\Gamma^m_{ij}\boldsymbol\varepsilon_m$:
\begin{align*}
\partial_jg_{ik}
&=\partial_j(\boldsymbol\varepsilon_i\cdot\boldsymbol\varepsilon_k)\\
&=(\Gamma^m_{ij}\boldsymbol\varepsilon_m)
      \cdot\boldsymbol\varepsilon_k
 +\boldsymbol\varepsilon_i\cdot
      (\Gamma^m_{kj}\boldsymbol\varepsilon_m)\\
&=\Gamma^m_{ij}g_{mk}+\Gamma^m_{kj}g_{im}.
\end{align*}
Use this identity in the product rule:
\begin{align*}
V_{i;j}
&=\partial_j(g_{ik}V^k)-V_\ell\Gamma^\ell_{ij}\\
&=g_{ik}\partial_jV^k+(\partial_jg_{ik})V^k
       -g_{\ell k}V^k\Gamma^\ell_{ij}\\
&=g_{ik}\partial_jV^k
  +\Gamma^m_{ij}g_{mk}V^k
  +\Gamma^m_{kj}g_{im}V^k
  -g_{mk}V^k\Gamma^m_{ij}\\
&=g_{ik}\partial_jV^k+g_{im}\Gamma^m_{kj}V^k\\
&=g_{ik}\bigl(\partial_jV^k+\Gamma^k_{mj}V^m\bigr)
 =\boxed{g_{ik}V^k{}_{;j}}.
\end{align*}
In the second line $V_\ell=g_{\ell k}V^k$; in the third the dummy
$\ell$ of the last term was renamed $m$. The two terms containing
$\Gamma^m_{ij}g_{mk}V^k$ then cancel exactly; only the dummy-index
relabeling $m\leftrightarrow k$ is used in the last line.

\Needspace{5\baselineskip}
\par\bigskip
\noindent\textbf{4.3.8}\quad From the circular cylindrical metric tensor
$g_{ij}$, calculate the $\Gamma_{ij}^{k}$ for circular cylindrical
coordinates.

\emph{Note.}~There are only three nonvanishing $\Gamma$.
\solution
For $q^i=(\rho,\phi,z)$, Exercise 4.3.5 gives
$g_{ij}=\operatorname{diag}(1,\rho^2,1)$ and
$g^{ij}=\operatorname{diag}(1,\rho^{-2},1)$. Use
\[
\Gamma^k_{ij}=\frac12g^{k\ell}
 (\partial_i g_{\ell j}+\partial_jg_{\ell i}-\partial_\ell g_{ij}).
\]
All metric derivatives vanish except $\partial_\rho g_{\phi\phi}=2\rho$.
For an upper $\rho$ index the inverse metric selects $\ell=\rho$;
only the last derivative can survive, and only when $i=j=\phi$:
\[
\Gamma^\rho_{\phi\phi}
=\frac12g^{\rho\rho}(0+0-\partial_\rho g_{\phi\phi})
=-\rho.
\]
For an upper $\phi$ index the inverse metric selects $\ell=\phi$. The
last derivative, $\partial_\phi g_{ij}$, is zero; the first survives
only for $i=\rho$, $j=\phi$, and the second only for $i=\phi$,
$j=\rho$:
\[
\Gamma^\phi_{\rho\phi}
=\frac12g^{\phi\phi}(\partial_\rho g_{\phi\phi}+0-0)
=\frac1{2\rho^2}(2\rho)=\frac1\rho,
\qquad
\Gamma^\phi_{\phi\rho}=\frac1\rho.
\]
With an upper $z$ index, every term differentiates a constant metric
entry or differentiates with respect to $z$, so every term is zero.
There are no other nonzero possibilities. Thus the complete list is
\[
\boxed{\Gamma^\rho_{\phi\phi}=-\rho,\qquad
\Gamma^\phi_{\rho\phi}=\Gamma^\phi_{\phi\rho}=\rho^{-1}}.
\]
As a check against the definition
$\partial_j\boldsymbol\varepsilon_i=\Gamma^k_{ij}\boldsymbol\varepsilon_k$,
differentiate the basis vectors of Exercise 4.3.5:
$\partial_\phi\boldsymbol\varepsilon_\phi
=(-\rho\cos\phi,-\rho\sin\phi,0)=-\rho\,\boldsymbol\varepsilon_\rho$ and
$\partial_\phi\boldsymbol\varepsilon_\rho
=\partial_\rho\boldsymbol\varepsilon_\phi
=(-\sin\phi,\cos\phi,0)=\boldsymbol\varepsilon_\phi/\rho$, while all
other derivatives vanish.

\Needspace{5\baselineskip}
\par\bigskip
\noindent\textbf{4.3.9}\quad Using the $\Gamma_{ij}^{k}$ from Exercise 4.3.8,
write out the covariant derivatives $V^{i}_{;j}$ of a vector $\mathbf{V}$ in
circular cylindrical coordinates.
\solution
Use the contravariant derivative formula
\[
V^i{}_{;j}=\partial_jV^i+\Gamma^i_{mj}V^m.
\]
For $i=\rho$, the only nonzero connection term occurs at $m=j=\phi$
and equals $-\rho V^\phi$. For $i=\phi$, the terms are
$V^\phi/\rho$ when $j=\rho$ and $V^\rho/\rho$ when $j=\phi$.
There are no connection terms for $i=z$. Thus all nine derivatives are
\[
\boxed{
\bigl(V^i{}_{;j}\bigr)=
\begin{pmatrix}
\partial_\rho V^\rho&\partial_\phi V^\rho-\rho V^\phi&\partial_z V^\rho\\[2pt]
\partial_\rho V^\phi+V^\phi/\rho&
\partial_\phi V^\phi+V^\rho/\rho&\partial_z V^\phi\\[2pt]
\partial_\rho V^z&\partial_\phi V^z&\partial_z V^z
\end{pmatrix}}.
\]
Rows are labeled by $i=(\rho,\phi,z)$ and columns by
$j=(\rho,\phi,z)$. These are components in the coordinate basis, whose
angular vector has length $\rho$. If
$\mathbf V=V_{\hat\rho}\hat{\mathbf e}_\rho+
V_{\hat\phi}\hat{\mathbf e}_\phi+V_{\hat z}\hat{\mathbf e}_z$,
then
\[
V^\rho=V_{\hat\rho},\qquad V^\phi=V_{\hat\phi}/\rho,
\qquad V^z=V_{\hat z}.
\]
For example, the first row's angular entry becomes
$\partial_\phi V^\rho-\rho V^\phi=\partial_\phi V_{\hat\rho}-V_{\hat\phi}$.
Since $\partial_j\mathbf V=V^i{}_{;j}\boldsymbol\varepsilon_i$ (the
relation that defines $V^i{}_{;j}$) and
$\boldsymbol\varepsilon_\rho=\hat{\mathbf e}_\rho$, this entry is the
$\hat{\mathbf e}_\rho$ component of $\partial_\phi\mathbf V$. Direct
differentiation of the unit-vector expansion, with
$\partial_\phi\hat{\mathbf e}_\rho=\hat{\mathbf e}_\phi$,
$\partial_\phi\hat{\mathbf e}_\phi=-\hat{\mathbf e}_\rho$ and
$\partial_\phi\hat{\mathbf e}_z=0$, agrees:
\[
\partial_\phi\mathbf V
=(\partial_\phi V_{\hat\rho}-V_{\hat\phi})\hat{\mathbf e}_\rho
 +(\partial_\phi V_{\hat\phi}+V_{\hat\rho})\hat{\mathbf e}_\phi
 +(\partial_\phi V_{\hat z})\hat{\mathbf e}_z.
\]
Likewise, because $\boldsymbol\varepsilon_\phi=\rho\,\hat{\mathbf e}_\phi$,
the $\hat{\mathbf e}_\phi$ component is
$\rho V^\phi{}_{;\phi}=\rho\,\partial_\phi(V_{\hat\phi}/\rho)+V_{\hat\rho}
=\partial_\phi V_{\hat\phi}+V_{\hat\rho}$.

\Needspace{5\baselineskip}
\par\bigskip
\noindent\textbf{4.3.10}\quad Show that for the metric tensor
$g_{ij;k}=g^{ij}_{;k}=0$.
\solution
The covariant derivative acts on each covariant index with a
minus connection term:
\[
g_{ij;k}=\partial_kg_{ij}
 -\Gamma^m_{ik}g_{mj}-\Gamma^m_{jk}g_{im}.
\]
This rule follows as in Exercise 4.3.6. Write a second-rank tensor as
$T_{ij}\boldsymbol\varepsilon^i\boldsymbol\varepsilon^j$, a sum of
direct products of reciprocal basis vectors, and differentiate by the
product rule, using
$\partial_k\boldsymbol\varepsilon^i=-\Gamma^i_{mk}\boldsymbol\varepsilon^m$
for each basis vector:
\begin{align*}
\partial_k(T_{ij}\boldsymbol\varepsilon^i\boldsymbol\varepsilon^j)
&=(\partial_kT_{ij})\boldsymbol\varepsilon^i\boldsymbol\varepsilon^j
 -T_{ij}\Gamma^i_{mk}\boldsymbol\varepsilon^m\boldsymbol\varepsilon^j
 -T_{ij}\Gamma^j_{mk}\boldsymbol\varepsilon^i\boldsymbol\varepsilon^m\\
&=(\partial_kT_{ij}-\Gamma^m_{ik}T_{mj}-\Gamma^m_{jk}T_{im})
 \boldsymbol\varepsilon^i\boldsymbol\varepsilon^j,
\end{align*}
after relabeling $i\leftrightarrow m$ in the second term and
$j\leftrightarrow m$ in the third; the coefficient is $T_{ij;k}$. With
contravariant components, $T^{ij}\boldsymbol\varepsilon_i
\boldsymbol\varepsilon_j$, and
$\partial_k\boldsymbol\varepsilon_i=\Gamma^m_{ik}\boldsymbol\varepsilon_m$
instead, each upper index gives a plus term.

Differentiate $g_{ij}=\boldsymbol\varepsilon_i\cdot
\boldsymbol\varepsilon_j$ as in Exercise 4.3.7 to obtain
\[
\partial_kg_{ij}=\Gamma^m_{ik}g_{mj}+\Gamma^m_{jk}g_{im}.
\]
Substitution cancels both terms, so $\boxed{g_{ij;k}=0}$.

For the inverse metric, differentiate $g^{ia}g_{ab}=\delta^i_b$:
\[
(\partial_kg^{ia})g_{ab}+g^{ia}\partial_kg_{ab}=0.
\]
Multiply by $g^{bj}$ and sum $b$. Since $g_{ab}g^{bj}=\delta^j_a$, the
first term becomes $\partial_kg^{ij}$, which is thereby isolated:
\begin{align*}
\partial_kg^{ij}
&=-g^{ia}(\partial_kg_{ab})g^{bj}\\
&=-g^{ia}(\Gamma^m_{ak}g_{mb}+\Gamma^m_{bk}g_{am})g^{bj}\\
&=-g^{ia}\Gamma^j_{ak}-\Gamma^i_{bk}g^{bj}.
\end{align*}
The last line uses $g_{mb}g^{bj}=\delta^j_m$ and
$g^{ia}g_{am}=\delta^i_m$. As shown above, the contravariant derivative
has a plus connection term for each upper index. Renaming the dummies
$a,b$ as $m$,
\begin{align*}
g^{ij}{}_{;k}
&=\partial_kg^{ij}+\Gamma^i_{mk}g^{mj}+\Gamma^j_{mk}g^{im}\\
&=-g^{im}\Gamma^j_{mk}-\Gamma^i_{mk}g^{mj}
  +\Gamma^i_{mk}g^{mj}+\Gamma^j_{mk}g^{im}=0.
\end{align*}
Both cancellations rely on the metric-compatible Christoffel
connection used in this section.

\Needspace{5\baselineskip}
\par\bigskip
\noindent\textbf{4.3.11}\quad Starting with the divergence in tensor notation,
Eq.~(4.70), develop the divergence of a vector in spherical polar coordinates,
Eq.~(3.157).
\solution
The tensor divergence is
\[
V^i{}_{;i}=\partial_iV^i+\Gamma^i_{mi}V^m
=\frac1{\sqrt g}\partial_i(\sqrt g\,V^i),
\qquad g=\det(g_{ij}).
\]
For completeness, the contraction in the middle expression satisfies
\[
\Gamma^i_{mi}=\frac12g^{i\ell}
 (\partial_mg_{\ell i}+\partial_i g_{\ell m}-\partial_\ell g_{mi})
=\frac12g^{i\ell}\partial_mg_{i\ell}
=\partial_m\ln\sqrt g.
\]
The last two derivatives in parentheses cancel after interchanging the
dummy indices $i,\ell$ in one of them and using the symmetry of
$g^{i\ell}$ and $g_{\ell m}$. The last equality is the determinant
derivative $\partial_mg=g\,g^{i\ell}\partial_mg_{\ell i}$: by the chain
rule, $\partial_mg$ is the sum over entries of
$(\partial g/\partial g_{\ell i})\,\partial_mg_{\ell i}$, and
$\partial g/\partial g_{\ell i}$ is the cofactor of $g_{\ell i}$, which
equals $g\,g^{i\ell}$ by the cofactor formula for the inverse matrix.
The product rule then gives the first displayed formula:
\[
\frac1{\sqrt g}\partial_i(\sqrt g\,V^i)
=\partial_iV^i+V^m\partial_m\ln\sqrt g
=\partial_iV^i+\Gamma^i_{mi}V^m.
\]

In spherical polar coordinates $q^i=(r,\theta,\phi)$, differentiating
$\mathbf r=(r\sin\theta\cos\phi,r\sin\theta\sin\phi,r\cos\theta)$ gives
the coordinate basis
\begin{align*}
\boldsymbol\varepsilon_r
&=(\sin\theta\cos\phi,\sin\theta\sin\phi,\cos\theta),\\
\boldsymbol\varepsilon_\theta
&=r(\cos\theta\cos\phi,\cos\theta\sin\phi,-\sin\theta),\\
\boldsymbol\varepsilon_\phi
&=r\sin\theta(-\sin\phi,\cos\phi,0).
\end{align*}
Their pairwise scalar products vanish, and their squared lengths are
$1$, $r^2$, $r^2\sin^2\theta$. Hence
\[
g_{ij}=\operatorname{diag}(1,r^2,r^2\sin^2\theta),\qquad
g=r^4\sin^2\theta,\qquad
\sqrt g=r^2\sin\theta
\]
on the usual regular chart $r>0$, $0<\theta<\pi$, where $\sin\theta>0$.
Let $V_r,V_\theta,V_\phi$ denote components along the \emph{unit} vectors.
The coordinate basis vectors have lengths $\sqrt{g_{ii}}=1,r,r\sin\theta$,
so comparing $\mathbf V=V^i\boldsymbol\varepsilon_i$ with the unit-vector
expansion gives
\[
V^r=V_r,\qquad V^\theta=\frac{V_\theta}{r},\qquad
V^\phi=\frac{V_\phi}{r\sin\theta}.
\]
Thus $\sqrt g\,V^r=r^2\sin\theta\,V_r$, $\sqrt g\,V^\theta=r\sin\theta\,V_\theta$,
and $\sqrt g\,V^\phi=rV_\phi$. Insert these quantities before differentiating:
\begin{align*}
\nabla\cdot\mathbf V
&=\frac1{r^2\sin\theta}
 \left[\partial_r(r^2\sin\theta V_r)
       +\partial_\theta(r\sin\theta V_\theta)
       +\partial_\phi(rV_\phi)\right]\\
&=\boxed{\frac1{r^2}\frac{\partial(r^2V_r)}{\partial r}
 +\frac1{r\sin\theta}\frac{\partial(\sin\theta V_\theta)}{\partial\theta}
 +\frac1{r\sin\theta}\frac{\partial V_\phi}{\partial\phi}}.
\end{align*}
Expanding the products gives the equivalent form
$\partial_rV_r+2V_r/r+(\partial_\theta V_\theta)/r+
(\cot\theta)V_\theta/r+(\partial_\phi V_\phi)/(r\sin\theta)$.

\Needspace{5\baselineskip}
\par\bigskip
\noindent\textbf{4.3.12}\quad The covariant vector $A_{i}$ is the gradient of
a \erratum{twice continuously differentiable} scalar. Show that the difference of covariant derivatives
$A_{i;j}-A_{j;i}$ vanishes.
\solution
Write $A_i=\partial_i\Phi$ for the scalar potential. Then
\[
A_{i;j}=\partial_j\partial_i\Phi-\Gamma^k_{ij}\partial_k\Phi,
\qquad
A_{j;i}=\partial_i\partial_j\Phi-\Gamma^k_{ji}\partial_k\Phi.
\]
Subtract:
\[
A_{i;j}-A_{j;i}
=(\partial_j\partial_i-\partial_i\partial_j)\Phi
 -(\Gamma^k_{ij}-\Gamma^k_{ji})\partial_k\Phi=0.
\]
The first term vanishes for a twice continuously differentiable scalar;
the second vanishes by Exercise 4.3.4. Thus the covariant Hessian of a
scalar is symmetric for this connection.

\Needspace{9\baselineskip}
\subsection*{4.4 Jacobians}

\Needspace{5\baselineskip}
\par\bigskip
\noindent\textbf{4.4.1}\quad Assuming the functions $u$ and $v$ to be
\erratum{continuously differentiable, and interpreting part (a) locally
where at least one of $\nabla u,\nabla v$ is nonzero},
\begin{enumerate}
\item[(a)] Show that a necessary and sufficient condition that $u(x,y,z)$ and
  $v(x,y,z)$ are related by \erratum{a regular $C^1$ relation $f(u,v)=0$, with
$(f_u,f_v)\ne(0,0)$}, is that
  $(\nabla u)\times(\nabla v)=0$;
\item[(b)] If $u=u(x,y)$ and $v=v(x,y)$, show that the condition
  $(\nabla u)\times(\nabla v)=0$ leads to the 2-D Jacobian
  \[
  J=\frac{\partial(u,v)}{\partial(x,y)}
  =\begin{vmatrix}\dfrac{\partial u}{\partial x}&\dfrac{\partial u}
  {\partial y}\\[2ex]\dfrac{\partial v}{\partial x}&\dfrac{\partial v}
  {\partial y}\end{vmatrix}=0.
  \]
\end{enumerate}
\solution
\erratum{The usual equivalence is local and needs regularity.
A meaningful relation is a $C^1$ function $f$ with
$(f_u,f_v)\ne(0,0)$ on the image near the point in question; otherwise
$f\equiv0$ would relate every pair $u,v$, and necessity would fail:
$u=x$, $v=y$ satisfy $f(u,v)=0$, yet
$\nabla u\times\nabla v=\hat{\mathbf e}_z\ne0$. For the converse, work where at
least one of $\nabla u,\nabla v$ is nonzero. Mere differentiability,
without these qualifications, does not establish the printed global
``necessary and sufficient'' statement.}

(a) First suppose $f(u(x,y,z),v(x,y,z))=0$ is a regular relation.
It holds identically near the point, so its partial derivatives vanish;
by the chain rule, $\partial_x f(u,v)=f_uu_x+f_vv_x=0$, and likewise
for $y$ and $z$. Together,
\[
f_u\nabla u+f_v\nabla v=0.
\]
If $f_v\ne0$, then $\nabla v=-(f_u/f_v)\nabla u$, and their cross
product is zero. If $f_v=0$, regularity forces $f_u\ne0$, whence
$\nabla u=0$ and the cross product again vanishes. This proves necessity.

Conversely suppose $\nabla u\times\nabla v=0$ in a neighborhood and
$\nabla u\ne0$ at the point under consideration. Relabel $x,y,z$ if
necessary so that $u_x\ne0$ there. Continuity lets us shrink the
neighborhood so that $u_x\ne0$ throughout. The coordinate change
\[
(s,\eta,\zeta)=(u(x,y,z),y,z)
\]
has Jacobian
\[
\frac{\partial(s,\eta,\zeta)}{\partial(x,y,z)}
=\begin{vmatrix}u_x&u_y&u_z\\0&1&0\\0&0&1\end{vmatrix}=u_x\ne0,
\]
so the inverse function theorem makes it locally invertible. The
vanishing cross product says that $\nabla v=\lambda\nabla u$, with
$\lambda=v_x/u_x$: its $z$- and $y$-components,
$u_xv_y-u_yv_x=0$ and $u_zv_x-u_xv_z=0$, give $v_y=\lambda u_y$ and
$v_z=\lambda u_z$, while $v_x=\lambda u_x$ by the choice of $\lambda$.
In differential notation this is $dv=\lambda\,du$. Let
$\widetilde v(s,\eta,\zeta)$ denote $v$ expressed in the new
coordinates. Since $s=u$, we have $du=ds$, and the relation reads
\[
d\widetilde v
=\frac{\partial\widetilde v}{\partial s}\,ds
 +\frac{\partial\widetilde v}{\partial\eta}\,d\eta
 +\frac{\partial\widetilde v}{\partial\zeta}\,d\zeta
=\lambda\,ds.
\]
The differentials $ds,d\eta,d\zeta$ of a coordinate system are linearly
independent, so the coefficients agree: $\widetilde v_s=\lambda$ and
$\widetilde v_\eta=\widetilde v_\zeta=0$. Let $(s_0,\eta_0,\zeta_0)$ be
the image of the point under consideration. On a sufficiently small
coordinate box $|s-s_0|<\delta$, $|\eta-\eta_0|<\delta$,
$|\zeta-\zeta_0|<\delta$, the mean value theorem applied in $\eta$ and
then in $\zeta$ gives
$\widetilde v(s,\eta,\zeta)=\widetilde v(s,\eta_0,\zeta_0)\equiv F(s)$,
a $C^1$ function of $s$ alone, because $\widetilde v$ is $v$ composed
with the inverse coordinate map, and both are $C^1$ (the inverse
function theorem supplies a $C^1$ inverse). Therefore
\[
v=F(u),\qquad f(u,v)=v-F(u)=0,
\]
and this relation is regular because $f_v=1$. If instead only
$\nabla v$ is known to be nonzero, interchange $u$ and $v$ in the
argument.

The restriction at critical points is substantive. For $u=x^2$,
$v=x^3$, the gradients $(2x,0,0)$ and $(3x^2,0,0)$ are parallel, so the
cross product vanishes everywhere, and both gradients vanish at $x=0$.
Yet no regular $C^1$ relation describes the image near
$(u,v)=(0,0)$. If $f_v(0,0)\ne0$, the implicit function theorem would
make the zero set of $f$ near the origin a graph $v=F(u)$; this is
contradicted by $x$ and $-x$ producing the same $u$ and opposite $v$.
If $f_u(0,0)\ne0$, it would give a $C^1$ function with $u=G(v)$, that
is, $G(x^3)=x^2$, so $G(v)=|v|^{2/3}$, whose difference quotient
$|v|^{2/3}/v$ at zero is unbounded. The singular relation
$v^2-u^3=0$, with $f_u=-3u^2$ and $f_v=2v$ both zero at the origin,
does exist. So at a critical point a vanishing cross product may yield
only a singular relation, which is why the converse in part (a) is
restricted to points where $\nabla u$ or $\nabla v$ is nonzero.

(b) If the functions are independent of $z$, write their gradients as
$(u_x,u_y,0)$ and $(v_x,v_y,0)$. Expanding the cross product gives
\[
\nabla u\times\nabla v
=\bigl(0,0,u_xv_y-u_yv_x\bigr)
=\hat{\mathbf e}_z\,
  \begin{vmatrix}u_x&u_y\\v_x&v_y\end{vmatrix}.
\]
Thus its vanishing is exactly
$\boxed{\partial(u,v)/\partial(x,y)=0}$. This algebraic identity
itself does not require the regular-point hypothesis in part (a).

\Needspace{5\baselineskip}
\par\bigskip
\noindent\textbf{4.4.2}\quad A 2-D orthogonal system is described by the
coordinates $q_{1}$ and $q_{2}$. Show that the Jacobian $J$ satisfies the
equation
\[
J\equiv\frac{\partial(x,y)}{\partial(q_{1},q_{2})}
\equiv\frac{\partial x}{\partial q_{1}}\frac{\partial y}{\partial q_{2}}
-\frac{\partial x}{\partial q_{2}}\frac{\partial y}{\partial q_{1}}=\erratum{\pm h_{1}h_{2}}.
\]
\emph{Hint.}~It's easier to work with the square of each side of this
equation.
\solution
\erratum{The orientation-independent statement is
$|J|=h_1h_2$. The textbook prints $J=h_1h_2$, whose positive sign
assumes an orientation-preserving coordinate chart; the orthogonal
chart $x=q_1$, $y=-q_2$ checked below has $J=-h_1h_2$.}

At a given point set
\[
a=x_{q_1},\quad b=y_{q_1},\quad c=x_{q_2},\quad d=y_{q_2}.
\]
The coordinate basis vectors are
$\partial\mathbf r/\partial q_1=(a,b)$ and
$\partial\mathbf r/\partial q_2=(c,d)$, and the scale factors are their
lengths, $h_i=|\partial\mathbf r/\partial q_i|$ (from
$ds^2=h_1^2dq_1^2+h_2^2dq_2^2$). Hence
\[
h_1^2=a^2+b^2,\qquad h_2^2=c^2+d^2,\qquad ac+bd=0.
\]
The last equality is precisely orthogonality,
$(\partial\mathbf r/\partial q_1)\cdot(\partial\mathbf r/\partial q_2)=0$.
Now expand the squares:
\begin{align*}
h_1^2h_2^2-(ac+bd)^2
&=(a^2+b^2)(c^2+d^2)-(a^2c^2+2abcd+b^2d^2)\\
&=a^2d^2+b^2c^2-2abcd\\
&=(ad-bc)^2=J^2.
\end{align*}
Because $ac+bd=0$, we obtain $J^2=h_1^2h_2^2$. Scale factors are
positive on a regular chart, hence $|J|=h_1h_2$. Positive orientation
selects $J=+h_1h_2$; negative orientation selects $J=-h_1h_2$.
For an explicit check, $x=q_1$, $y=-q_2$ is orthogonal with
$h_1=h_2=1$ but $J=-1$, so orthogonality alone cannot fix the sign.

\Needspace{5\baselineskip}
\par\bigskip
\noindent\textbf{4.4.3}\quad For the transformation $u=x+y$, $v=x/y$, with
$x\ge0$ and \erratum{$y>0$}, find the Jacobian
$\dfrac{\partial(x,y)}{\partial(u,v)}$
\begin{enumerate}
\item[(a)] By direct computation,
\item[(b)] By first computing $J^{-1}$.
\end{enumerate}
\solution
\erratum{The ratio $v=x/y$ requires $y>0$ on the stated
nonnegative quadrant; the boundary $y=0$ is excluded.} Then $u>0$,
$v\ge0$, and the inverse transformation follows directly from
$x=vy$:
\[
u=x+y=(v+1)y,\qquad
\boxed{y=\frac{u}{v+1},\qquad x=\frac{uv}{v+1}}.
\]
(a) Differentiate the inverse formulas; for $x_v$ the quotient rule
gives $x_v=u[(v+1)-v]/(v+1)^2$:
\[
x_u=\frac{v}{v+1},\qquad x_v=\frac{u}{(v+1)^2},\qquad
y_u=\frac1{v+1},\qquad y_v=-\frac{u}{(v+1)^2}.
\]
Their determinant is
\begin{align*}
J&=x_uy_v-x_vy_u
 =\frac{v}{v+1}\left(-\frac{u}{(v+1)^2}\right)
  -\frac{u}{(v+1)^2}\cdot\frac1{v+1}\\
 &=-\frac{uv}{(v+1)^3}-\frac{u}{(v+1)^3}
 =-\frac{u(v+1)}{(v+1)^3}
 =\boxed{-\frac{u}{(v+1)^2}}.
\end{align*}
(b) For the forward transformation,
\[
u_x=1,\quad u_y=1,\quad v_x=\frac1y,\quad v_y=-\frac{x}{y^2}.
\]
Thus
\[
J^{-1}=\frac{\partial(u,v)}{\partial(x,y)}
=\begin{vmatrix}1&1\\1/y&-x/y^2\end{vmatrix}
=-\frac{x}{y^2}-\frac1y
=-\frac{x+y}{y^2}=-\frac{u}{y^2}.
\]
It is nonzero on the admissible domain. The chain rule makes the
Jacobian matrices of a map and its inverse mutually inverse, and the
determinant of a product of matrices is the product of their
determinants, so
$\bigl[\partial(x,y)/\partial(u,v)\bigr]\bigl[\partial(u,v)/\partial(x,y)\bigr]=1$,
and with $y=u/(v+1)$,
\[
J=-\frac{y^2}{u}
=-\frac1u\left(\frac{u}{v+1}\right)^2
=-\frac{u}{(v+1)^2},
\]
in agreement with (a). The negative sign records orientation reversal;
an ordinary unoriented area integral uses $|J|=u/(v+1)^2$.

\Needspace{9\baselineskip}
\subsection*{4.5 Differential Forms}

\Needspace{5\baselineskip}
\par\bigskip
\noindent\textbf{4.5.1}\quad Using the rules for the application of the Hodge
star operator, verify the results given in Eq.~(4.82) for its application to
all linearly independent differential forms in Minkowski space.
\erratum{In the rule on p.~236, count the negative metric factors in the
input form being starred, rather than in its complement.}
\solution
Use the ordered basis of 1-forms
\[
e^0=dt,\quad e^1=dx_1,\quad e^2=dx_2,\quad e^3=dx_3,
\qquad \eta=\operatorname{diag}(1,-1,-1,-1),
\]
and write $e^{a_1\cdots a_p}=e^{a_1}\wedge\cdots\wedge e^{a_p}$,
with $\mathrm{vol}=e^{0123}$. For an increasing index set $I$, let $J$
be its increasing complement, $m$ the number of spatial indices in $I$,
and $N$ the number of inversions in the concatenation $(I,J)$, that is,
of pairs of entries standing in decreasing order; for example, for
$I=\{1,3\}$ the concatenation $(1,3,0,2)$ has the three inversions
$(1,0)$, $(3,0)$, $(3,2)$, so $N=3$. Then
\[
e^I\wedge e^J=(-1)^N\mathrm{vol},\qquad
\langle e^I,e^I\rangle=(-1)^m{:}
\]
the first because sorting $(I,J)$ into the order $0123$ takes $N$
interchanges of adjacent factors, each of which changes the sign (each
interchange of an adjacent pair standing in decreasing order removes
exactly one inversion); the second because the metric induces on
$p$-forms the inner product
$\langle e^K,e^I\rangle=\det\bigl[\eta^{ab}\bigr]_{a\in K,\,b\in I}$,
which for $K=I$ is the product of the diagonal metric entries
$\eta^{aa}$, $a\in I$, $m$ of which equal $-1$. For $K\ne I$ this
determinant vanishes: since $\eta$ is diagonal, the row of an index
$a\in K$ not in $I$ is zero.
The Hodge star of a $p$-form $\beta$ is fixed by the standard defining
relation
$\alpha\wedge *\beta=\langle\alpha,\beta\rangle\mathrm{vol}$ for every
$p$-form $\alpha$. Take $\beta=e^I$. For $\alpha=e^K$ with $K\ne I$ the
right side is zero by this orthogonality, and $e^K$ has a nonzero wedge
product only with the basis form of its own complement; hence $*e^I$
has no component along the complement of any $K\ne I$, and
$*e^I=s\,e^J$ for some number $s$.
Then $\alpha=e^I$ gives $s(-1)^N=(-1)^m$, so
\[
\boxed{*e^I=(-1)^{N+m}e^J}.
\]
The sign combines the parity $(-1)^N$ of the permutation that orders
$e^I\wedge e^J$ with a factor $-1$ for each negative metric entry of
the form $e^I$ being starred.

\erratum{The rule in the first paragraph of textbook p.~236 says to
count negative metric entries in the complementary form $\omega'$.
That is a typographical error: count them in the form $\omega$ being
starred. The worked example and Eq.~(4.82) use this corrected rule:
for $\omega=1$ the complement $\omega'=\mathrm{vol}$ has three negative
metric factors, so the printed rule would give $*1=-\mathrm{vol}$,
whereas Eq.~(4.82) has $*1=+\mathrm{vol}$. Since $\omega$ and $\omega'$
together contain all three spatial differentials, the two counts always
differ in parity, so the printed rule would reverse the sign of every
result in Eq.~(4.82).}
For instance, $*1=\mathrm{vol}$ has $m=0$, while
$*\mathrm{vol}=-1$ has $m=3$.

\Needspace{22\baselineskip}
Here is the explicit calculation for every basis form. The complement
column is written before its final sign is applied.
\[
\begin{array}{c|c|r|r|c}
\text{input}&\text{complement}&N&m&\text{Hodge star}\\ \hline
1&e^{0123}&0&0&e^{0123}\\
e^0&e^{123}&0&0&e^{123}\\
e^1&e^{023}&1&1&e^{023}\\
e^2&e^{013}&2&1&-e^{013}\\
e^3&e^{012}&3&1&e^{012}\\ \hline
e^{01}&e^{23}&0&1&-e^{23}\\
e^{02}&e^{13}&1&1&e^{13}\\
e^{03}&e^{12}&2&1&-e^{12}\\
e^{12}&e^{03}&2&2&e^{03}\\
e^{13}&e^{02}&3&2&-e^{02}\\
e^{23}&e^{01}&4&2&e^{01}\\ \hline
e^{012}&e^3&0&2&e^3\\
e^{013}&e^2&1&2&-e^2\\
e^{023}&e^1&2&2&e^1\\
e^{123}&e^0&3&3&e^0\\
e^{0123}&1&0&3&-1
\end{array}
\]
The table includes $1+4+6+4+1=16$ independent forms. Take $(i,j,k)$
to be $(1,2,3)$, $(2,3,1)$, or $(3,1,2)$. Rewriting the table in cyclic order
absorbs the displayed minus signs associated with reordering spatial
indices. These occur in the rows $e^2$, $e^{13}$, $e^{013}$; using
$dx_1\wedge dx_3=-dx_3\wedge dx_1$ (also inside the starred forms),
\begin{gather*}
*dx_2=-dt\wedge dx_1\wedge dx_3=dt\wedge dx_3\wedge dx_1,\qquad
*(dx_3\wedge dx_1)=-{*}(dx_1\wedge dx_3)=dt\wedge dx_2,\\
*(dt\wedge dx_3\wedge dx_1)=-{*}(dt\wedge dx_1\wedge dx_3)=dx_2,
\end{gather*}
while the row $e^{02}$ gives
$*(dt\wedge dx_2)=dx_1\wedge dx_3=-dx_3\wedge dx_1$.
The table thus gives exactly
\begin{gather*}
*1=dt\wedge dx_1\wedge dx_2\wedge dx_3,\qquad
*dt=dx_1\wedge dx_2\wedge dx_3,\\
*dx_i=dt\wedge dx_j\wedge dx_k,\\
*(dx_j\wedge dx_k)=dt\wedge dx_i,\qquad
*(dt\wedge dx_i)=-dx_j\wedge dx_k,\\
*(dx_1\wedge dx_2\wedge dx_3)=dt,\qquad
*(dt\wedge dx_i\wedge dx_j)=dx_k,\\
*(dt\wedge dx_1\wedge dx_2\wedge dx_3)=-1.
\end{gather*}
As an independent sign check, apply the boxed rule twice to a $p$-form:
$*e^I=(-1)^{N_I+m_I}e^J$ and $*e^J=(-1)^{N_J+m_J}e^I$, where $N_J$
counts the inversions of $(J,I)$. Since $I$ and $J$ are both listed in
increasing order, no two indices from the same set form an inversion,
while each of the $p(4-p)$ pairs made of one index from $I$ and one
from $J$ is an inversion in exactly one of the orders $(I,J)$ and
$(J,I)$; so $N_I+N_J=p(4-p)$. Also $m_I+m_J=3$, the total number of
spatial indices. Hence
\[
{*}{*}\omega_p=(-1)^{p(4-p)+3}\omega_p,
\]
which is negative for even $p$ and positive for odd $p$, in agreement
with every pair of rows in the table.

\Needspace{5\baselineskip}
\par\bigskip
\noindent\textbf{4.5.2}\quad If the force field is constant and moving a
particle from the origin to $(3,0,0)$ requires $a$ units of work, from
$(-1,-1,0)$ to $(-1,1,0)$ takes $b$ units of work, and from $(0,0,4)$ to
$(0,0,5)$ $c$ units of work, find the 1-form of the work.
\solution
Let $\mathbf F=(F_x,F_y,F_z)$ be the constant force that does the quoted
work (if the work is done against the field, $\mathbf F$ is minus the
field's force). Its work 1-form is
$\omega=F_x\,dx+F_y\,dy+F_z\,dz$, and its work between two points is
$\mathbf F\cdot\Delta\mathbf r$, because each force component can be
taken outside the line integral. The three given displacements are
\[
(3,0,0),\qquad (0,2,0),\qquad (0,0,1).
\]
They yield $3F_x=a$, $2F_y=b$, and $F_z=c$. Therefore
\[
\boxed{\omega=\frac a3\,dx+\frac b2\,dy+c\,dz}.
\]
For a direct check, integrating along the three paths multiplies these
coefficients by $3,2,1$, respectively, returning $a,b,c$. The form is
also exact:
$\omega=d(ax/3+by/2+cz)$, so the result does not depend on the shapes
of the paths.

\Needspace{9\baselineskip}
\subsection*{4.6 Differentiating Forms}

\Needspace{5\baselineskip}
\par\bigskip
\noindent\textbf{4.6.1}\quad Given the two 1-forms $\omega_{1}=x\,dy+y\,dx$
and $\omega_{2}=x\,dy-y\,dx$, calculate
\begin{enumerate}
\item[(a)] $d\omega_{1}$,
\item[(b)] $d\omega_{2}$.
\item[(c)] For each of your answers to (a) or (b) that is nonzero, apply the
  operator $d$ a second time and verify that $d(d\omega_{i})=0$.
\end{enumerate}
\solution
Use $d(f\alpha)=df\wedge\alpha+f\,d\alpha$ and
$d(dx)=d(dy)=0$.

(a) For the first form,
\begin{align*}
d\omega_1
&=d(x\,dy)+d(y\,dx)\\
&=dx\wedge dy+x\,d(dy)+dy\wedge dx+y\,d(dx)\\
&=dx\wedge dy-dx\wedge dy=\boxed{0}.
\end{align*}
Indeed $\omega_1=d(xy)$.

(b) Changing the sign between the terms, and using the two derivatives
$d(x\,dy)=dx\wedge dy$ and $d(y\,dx)=dy\wedge dx$ found in (a), gives
\begin{align*}
d\omega_2
&=d(x\,dy)-d(y\,dx)\\
&=dx\wedge dy-dy\wedge dx
=\boxed{2\,dx\wedge dy}.
\end{align*}

(c) The nonzero result is a constant-coefficient 2-form. Applying $d$
again, and using the product rule
$d(\alpha\wedge\beta)=d\alpha\wedge\beta+(-1)^p\,\alpha\wedge d\beta$
for a $p$-form $\alpha$ (here $\alpha=dx$ with $p=1$, hence the minus
sign),
\[
d(d\omega_2)
=2\,d(dx\wedge dy)
=2\{d(dx)\wedge dy-dx\wedge d(dy)\}=0.
\]
The already-zero result for $\omega_1$ also has zero exterior derivative.

\Needspace{5\baselineskip}
\par\bigskip
\noindent\textbf{4.6.2}\quad Apply the operator $d$ twice to
$\omega_{3}=xy\,dz+xz\,dy-yz\,dx$. Verify that the second application of $d$
yields a zero result.
\solution
First differentiate each coefficient and keep every wedge term:
\begin{align*}
d\omega_3
&=d(xy)\wedge dz+d(xz)\wedge dy-d(yz)\wedge dx\\
&=(y\,dx+x\,dy)\wedge dz
 +(z\,dx+x\,dz)\wedge dy
 -(z\,dy+y\,dz)\wedge dx\\
&=y\,dx\wedge dz+x\,dy\wedge dz
 +z\,dx\wedge dy+x\,dz\wedge dy\\
&\hspace{1.5em}-z\,dy\wedge dx-y\,dz\wedge dx.
\end{align*}
The two terms with coefficient $x$ cancel. Reordering the last two
terms reverses their signs, so the terms with coefficients $y$ and $z$
each add rather than cancel:
\[
\boxed{d\omega_3=2y\,dx\wedge dz+2z\,dx\wedge dy}.
\]
Apply $d$ again. Derivatives of the basis differentials vanish, leaving
$d(2y)\wedge dx\wedge dz+d(2z)\wedge dx\wedge dy$ with $d(2y)=2\,dy$
and $d(2z)=2\,dz$:
\begin{align*}
d(d\omega_3)
&=2\,dy\wedge dx\wedge dz+2\,dz\wedge dx\wedge dy\\
&=-2\,dx\wedge dy\wedge dz+2\,dx\wedge dy\wedge dz=\boxed{0}.
\end{align*}
The first permutation needs one interchange; the second needs two,
which explains their opposite signs.

\Needspace{5\baselineskip}
\par\bigskip
\noindent\textbf{4.6.3}\quad For $\omega_{2}$ and $\omega_{3}$ the 1-forms
with these names in Exercises 4.6.1 and 4.6.2, evaluate
$d(\omega_{2}\wedge\omega_{3})$:
\begin{enumerate}
\item[(a)] By forming the exterior product and then differentiating, and
\item[(b)] Using the formula for differentiating a product of two forms.
\end{enumerate}
Verify that both approaches give the same result.
\solution
(a) Expand the product, keeping the order of differentials:
\begin{align*}
\omega_2\wedge\omega_3
&=(x\,dy-y\,dx)\wedge(xy\,dz+xz\,dy-yz\,dx)\\
&=x^2y\,dy\wedge dz+x^2z\,dy\wedge dy
 -xyz\,dy\wedge dx\\
&\hspace{1.5em}-xy^2\,dx\wedge dz-xyz\,dx\wedge dy
 +y^2z\,dx\wedge dx\\
&=x^2y\,dy\wedge dz-xy^2\,dx\wedge dz.
\end{align*}
The repeated-differential terms vanish, and the two $xyz$ terms cancel
because $dy\wedge dx=-dx\wedge dy$. Differentiating the surviving
coefficients, $d(x^2y)=2xy\,dx+x^2\,dy$ and $d(xy^2)=y^2\,dx+2xy\,dy$,
gives the following; in the third line the terms
$x^2\,dy\wedge dy\wedge dz$ and $-y^2\,dx\wedge dx\wedge dz$ are dropped
because they contain a repeated differential:
\begin{align*}
d(\omega_2\wedge\omega_3)
&=(2xy\,dx+x^2\,dy)\wedge dy\wedge dz\\
&\hspace{1.5em}-(y^2\,dx+2xy\,dy)\wedge dx\wedge dz\\
&=2xy\,dx\wedge dy\wedge dz-2xy\,dy\wedge dx\wedge dz\\
&=2xy\,dx\wedge dy\wedge dz+2xy\,dx\wedge dy\wedge dz
=\boxed{4xy\,dx\wedge dy\wedge dz}.
\end{align*}

(b) Since $\omega_2$ is a 1-form, the product rule quoted in
Exercise 4.6.1(c) has a minus sign:
\[
d(\omega_2\wedge\omega_3)
=d\omega_2\wedge\omega_3-\omega_2\wedge d\omega_3.
\]
Use the preceding exercises to compute the terms separately:
\begin{align*}
d\omega_2\wedge\omega_3
&=2\,dx\wedge dy\wedge(xy\,dz+xz\,dy-yz\,dx)\\
&=2xy\,dx\wedge dy\wedge dz,\\
\omega_2\wedge d\omega_3
&=(x\,dy-y\,dx)\wedge(2y\,dx\wedge dz+2z\,dx\wedge dy)\\
&=2xy\,dy\wedge dx\wedge dz
=-2xy\,dx\wedge dy\wedge dz.
\end{align*}
All terms not displayed in the simplified lines contain a repeated
$dx$ or $dy$ and are zero. Subtracting the second result from the first
gives $(2xy-(-2xy))\,dx\wedge dy\wedge dz=4xy\,dx\wedge dy\wedge dz$,
exactly as in (a).

\Needspace{9\baselineskip}
\subsection*{4.7 Integrating Forms}

\Needspace{5\baselineskip}
\par\bigskip
\noindent\textbf{4.7.1}\quad Use differential-forms relations to transform the
integral $A(x,y,z)\,dx\wedge dy\wedge dz$ to the equivalent expression in
$du\wedge dv\wedge dw$, where $u,v,w$ is an \erratum{invertible} linear transformation of $x,y,z$,
and thereby find the determinant that can be identified as the Jacobian of the
transformation.
\solution
\erratum{To serve as a change of three independent coordinates,
the linear transformation must be invertible.} Write its inverse as
\[
\begin{pmatrix}x\\y\\z\end{pmatrix}
=\begin{pmatrix}a_1&a_2&a_3\\b_1&b_2&b_3\\c_1&c_2&c_3\end{pmatrix}
 \begin{pmatrix}u\\v\\w\end{pmatrix}
\equiv L\begin{pmatrix}u\\v\\w\end{pmatrix}.
\]
The coefficients are constant, so
\[
dx=a_1\,du+a_2\,dv+a_3\,dw,\quad
dy=b_1\,du+b_2\,dv+b_3\,dw,\quad
dz=c_1\,du+c_2\,dv+c_3\,dw.
\]
Expand their exterior product. Any term containing two copies of a
differential vanishes; precisely six terms remain, one for each
ordering of $du,dv,dw$:
\begin{align*}
dx\wedge dy\wedge dz
={}&a_1b_2c_3\,du\wedge dv\wedge dw
 +a_1b_3c_2\,du\wedge dw\wedge dv\\
&+a_2b_1c_3\,dv\wedge du\wedge dw
 +a_2b_3c_1\,dv\wedge dw\wedge du\\
&+a_3b_1c_2\,dw\wedge du\wedge dv
 +a_3b_2c_1\,dw\wedge dv\wedge du.
\end{align*}
Each interchange of two factors changes the sign. The orderings
$du\wedge dw\wedge dv$, $dv\wedge du\wedge dw$, $dw\wedge dv\wedge du$
are one interchange away from $du\wedge dv\wedge dw$ and acquire a
factor $-1$; the cyclic shifts $dv\wedge dw\wedge du$ and
$dw\wedge du\wedge dv$ are two interchanges away and keep their sign.
Hence
\[
dx\wedge dy\wedge dz
=(a_1b_2c_3-a_1b_3c_2-a_2b_1c_3+a_2b_3c_1
     +a_3b_1c_2-a_3b_2c_1)\,du\wedge dv\wedge dw.
\]
The coefficient is
$a_1(b_2c_3-b_3c_2)-a_2(b_1c_3-b_3c_1)+a_3(b_1c_2-b_2c_1)$, the
first-row expansion of the determinant of $L$. Since $x_u=a_1$,
$x_v=a_2$, \dots, $z_w=c_3$, it is
\[
J=\det L
=\begin{vmatrix}x_u&x_v&x_w\\y_u&y_v&y_w\\z_u&z_v&z_w\end{vmatrix}
=\frac{\partial(x,y,z)}{\partial(u,v,w)}.
\]
Therefore the transformed 3-form is
\[
\boxed{A(x(u,v,w),y(u,v,w),z(u,v,w))\,J\,du\wedge dv\wedge dw}.
\]
Its integral is taken over the corresponding parameter region with the
induced orientation. Differential forms retain the sign of $J$;
a scalar integral written using positive, unoriented volume elements
instead uses $|J|$. If the forward transformation has matrix $M$,
then $L=M^{-1}$ and $J=1/\det M$.

\Needspace{5\baselineskip}
\par\bigskip
\noindent\textbf{4.7.2}\quad Write Oersted's law,
\[
\int_{\partial S}\mathbf{H}\cdot d\mathbf{r}
=\int_{S}\nabla\times\mathbf{H}\cdot d\mathbf{a}\sim I,
\]
in differential form notation.
\solution
Associate the 1-form
\[
h=H_x\,dx+H_y\,dy+H_z\,dz
\]
with the magnetic field. Its exterior derivative is
$dh=dH_x\wedge dx+dH_y\wedge dy+dH_z\wedge dz$. Expanding each
$dH_i=\partial_xH_i\,dx+\partial_yH_i\,dy+\partial_zH_i\,dz$, dropping
the terms with a repeated differential, and writing every pair in the
order $dy\wedge dz$, $dz\wedge dx$, or $dx\wedge dy$,
\begin{align*}
dH_x\wedge dx&=\partial_yH_x\,dy\wedge dx+\partial_zH_x\,dz\wedge dx
 =-\partial_yH_x\,dx\wedge dy+\partial_zH_x\,dz\wedge dx,\\
dH_y\wedge dy&=\partial_xH_y\,dx\wedge dy+\partial_zH_y\,dz\wedge dy
 =\partial_xH_y\,dx\wedge dy-\partial_zH_y\,dy\wedge dz,\\
dH_z\wedge dz&=\partial_xH_z\,dx\wedge dz+\partial_yH_z\,dy\wedge dz
 =-\partial_xH_z\,dz\wedge dx+\partial_yH_z\,dy\wedge dz.
\end{align*}
Grouping terms with the same ordered differential pair,
\begin{align*}
dh={}&(\partial_yH_z-\partial_zH_y)\,dy\wedge dz\\
&+(\partial_zH_x-\partial_xH_z)\,dz\wedge dx\\
&+(\partial_xH_y-\partial_yH_x)\,dx\wedge dy.
\end{align*}
These three coefficients are the components of $\nabla\times\mathbf H$.
To integrate over $S$, parametrize it as $\mathbf r(s,t)$, with $(s,t)$
ordered so that $\mathbf r_s\times\mathbf r_t$ points along the chosen
normal; then $d\mathbf a=\mathbf r_s\times\mathbf r_t\,ds\,dt$. On $S$,
\[
dy\wedge dz=(y_s\,ds+y_t\,dt)\wedge(z_s\,ds+z_t\,dt)
=(y_sz_t-y_tz_s)\,ds\wedge dt,
\]
and $y_sz_t-y_tz_s$ is the $x$-component of
$\mathbf r_s\times\mathbf r_t$; cyclically ($x\to y\to z\to x$),
$dz\wedge dx$ and $dx\wedge dy$ give its $y$- and $z$-components.
Hence a 2-form $P\,dy\wedge dz+Q\,dz\wedge dx+R\,dx\wedge dy$
integrates over $S$ to the flux $\int_S(P,Q,R)\cdot d\mathbf a$; in
particular $\int_Sdh=\int_S\nabla\times\mathbf H\cdot d\mathbf a$.
Along $\partial S$ the 1-form $h$ is $\mathbf H\cdot d\mathbf r$, so
$\int_{\partial S}h=\int_{\partial S}\mathbf H\cdot d\mathbf r$.
Stokes' theorem thus gives the requested notation:
\[
\boxed{\int_{\partial S}h=\int_Sdh\ \sim\ I}.
\]
The boundary orientation is the one induced by $S$. More explicitly,
if the current density is $\mathbf J$, its flux 2-form is
\[
j=J_x\,dy\wedge dz+J_y\,dz\wedge dx+J_z\,dx\wedge dy,
\qquad I=\int_S j.
\]
With the proportionality constant $\kappa$ appropriate to the chosen
units ($\kappa=1$ in SI units, where $\nabla\times\mathbf H=\mathbf J$;
$\kappa=4\pi/c$ in Gaussian units), the magnetostatic field equation is
$dh=\kappa j$, and the
integral relation reads $\int_{\partial S}h=\kappa\int_Sj=\kappa I$.
This preserves the unspecified proportionality convention in the
question.

\Needspace{5\baselineskip}
\par\bigskip
\noindent\textbf{4.7.3}\quad A 1-form $A\,dx+B\,dy$ is defined as closed if
$\dfrac{\partial A}{\partial y}=\dfrac{\partial B}{\partial x}$. It is called
exact if there is a function $f$ such that $\dfrac{\partial f}{\partial x}=A$
and $\dfrac{\partial f}{\partial y}=B$. Determine which of the following
1-forms are closed, or exact, and find the corresponding functions $f$ for
those that are exact.
\erratum{Use their natural real domains for the first four forms, and
state explicitly when a primitive is only local. For the fifth, take a
$C^1$ complex-valued coefficient on an open domain in $\mathbb C$ and give
the conditions for closedness and global exactness:}
\[
y\,dx+x\,dy,\quad\frac{y\,dx+x\,dy}{x^{2}+y^{2}},\quad
[\ln(xy)+1]\,dx+\frac{x}{y}\,dy,
\]
\[
-\frac{y\,dx}{x^{2}+y^{2}}+\frac{x\,dy}{x^{2}+y^{2}},\quad
f(z)\,dz\text{ with }z=x+iy.
\]
\solution
For a $C^1$ 1-form $\omega=A\,dx+B\,dy$, the terms $A_x\,dx\wedge dx$
and $B_y\,dy\wedge dy$ vanish, so
\[
d\omega=dA\wedge dx+dB\wedge dy
=A_y\,dy\wedge dx+B_x\,dx\wedge dy
=(B_x-A_y)\,dx\wedge dy.
\]
Thus closedness means $B_x=A_y$. A globally exact form $\omega=d\Phi$
is closed, because $A_y=\Phi_{xy}=\Phi_{yx}=B_x$ (equality of mixed
partial derivatives, Schwarz's theorem; $\Phi$ is $C^2$ since
its first derivatives $A,B$ are $C^1$), and it has zero integral around
every closed curve $\Gamma$, because $\int_\Gamma d\Phi$ is the
difference of the values of $\Phi$ at the endpoints of $\Gamma$, which
coincide. We use
$\Phi$ for a potential to distinguish it from the last question's $f$.

\erratum{Exactness depends on the domain. In particular the angular
form below is only locally exact on the punctured plane. For the last
form, closedness is equivalent to holomorphy of $f$, and a closed
$f\,dz$ is globally exact exactly when $\oint_\Gamma f\,dz=0$ for every
closed curve $\Gamma$ in the domain, which holds on every simply
connected domain but not in general ($dz/z$ on
$\mathbb C\setminus\{0\}$ is a counterexample).}

\emph{First form.} For $\omega=y\,dx+x\,dy$, $A_y=B_x=1$, so the
form is closed. Integrate $\Phi_x=y$ with respect to $x$:
$\Phi=xy+h(y)$. The other equation gives
$\Phi_y=x+h'(y)=x$, hence $h'=0$. Therefore
\[
\boxed{\omega=d(xy),\qquad\Phi=xy+C},
\]
and the form is globally exact on $\mathbb R^2$.

\emph{Second form.} Its domain excludes the origin. Set
$A=y/(x^2+y^2)$ and $B=x/(x^2+y^2)$. The quotient rule gives
\[
A_y=\frac{x^2+y^2-2y^2}{(x^2+y^2)^2}
=\frac{x^2-y^2}{(x^2+y^2)^2},\qquad
B_x=\frac{x^2+y^2-2x^2}{(x^2+y^2)^2}
=\frac{y^2-x^2}{(x^2+y^2)^2}.
\]
These derivatives differ except on the lines $y=\pm x$. Those lines
contain no nonempty open subset of the plane, and a form that is exact
on an open set is closed there. Consequently
\[
\boxed{\frac{y\,dx+x\,dy}{x^2+y^2}\text{ is neither closed nor exact}}
\]
on any nonempty open part of its natural domain.

\emph{Third form.} For real logarithms the domain is $xy>0$.
The two coefficients obey
\[
\partial_y[\ln(xy)+1]=\frac{x}{xy}=\frac1y,
\qquad \partial_x(x/y)=\frac1y,
\]
so the form is closed. Integration in $x$ gives
$\Phi=x\ln(xy)+h(y)$, because, by the product rule and
$\partial_x\ln(xy)=y/(xy)=1/x$,
$\partial_x[x\ln(xy)]=\ln(xy)+x\cdot(1/x)=\ln(xy)+1$. Then
$\Phi_y=x/y+h'(y)=x/y$, so $h$ is constant. Thus
\[
\boxed{[\ln(xy)+1]\,dx+\frac xy\,dy=d[x\ln(xy)]}.
\]
It is exact on each component of $xy>0$, and also on their union
with the displayed potential; arbitrary additive constants may be
chosen independently on the two components.

\emph{Fourth form.} Here $A=-y/(x^2+y^2)$ and
$B=x/(x^2+y^2)$, so
\[
A_y=\frac{-(x^2+y^2)+2y^2}{(x^2+y^2)^2}
=\frac{y^2-x^2}{(x^2+y^2)^2}=B_x,
\]
with $B_x$ as computed for the second form.
The form is closed everywhere on $\mathbb R^2\setminus\{0\}$.
On a region with a continuous angle coordinate $\theta$, use
$x=r\cos\theta$, $y=r\sin\theta$, so that
$dx=\cos\theta\,dr-r\sin\theta\,d\theta$ and
$dy=\sin\theta\,dr+r\cos\theta\,d\theta$. The $dr$ terms cancel and the
$d\theta$ terms combine through $\sin^2\theta+\cos^2\theta=1$:
\begin{align*}
-y\,dx+x\,dy
={}&-r\sin\theta(\cos\theta\,dr-r\sin\theta\,d\theta)\\
&+r\cos\theta(\sin\theta\,dr+r\cos\theta\,d\theta)
=r^2\,d\theta.
\end{align*}
Dividing by $x^2+y^2=r^2$, $\omega=d\theta$ locally. For example, on
$x>0$ a potential is
$\Phi=\arctan(y/x)+C$, since
\[
\Phi_x=\frac{-y/x^2}{1+y^2/x^2}=-\frac{y}{x^2+y^2},\qquad
\Phi_y=\frac{1/x}{1+y^2/x^2}=\frac{x}{x^2+y^2}.
\]
For a complete counterclockwise circle
$x=R\cos t$, $y=R\sin t$, $0\le t\le2\pi$, where
$dx=-R\sin t\,dt$ and $dy=R\cos t\,dt$,
\[
\oint\omega
=\int_0^{2\pi}\frac{R^2\sin^2t+R^2\cos^2t}{R^2}\,dt=2\pi.
\]
A single-valued potential would make this integral zero. Therefore
\[
\boxed{\omega\text{ is closed but not globally exact on the punctured plane}}.
\]
\erratum{The Instructor's Manual, p.~65, prints $d\tan(y/x)$;
the local primitive is $\arctan(y/x)$, and its assertion of exactness
needs the domain qualification just proved.}

\emph{Fifth form.} Write the complex-valued coefficient as $f(x,y)$
and use $dz=dx+i\,dy$. Thus $A=f$ and $B=if$, so
$B_x-A_y=if_x-f_y=2i\cdot\frac12(f_x+if_y)$, giving
\[
d(f\,dz)=(if_x-f_y)\,dx\wedge dy
=2i\frac{\partial f}{\partial\bar z}\,dx\wedge dy,
\qquad
\frac{\partial f}{\partial\bar z}=\frac12(f_x+if_y).
\]
For $C^1$ coefficients the form is closed exactly when
$f_y=if_x$. With $f=u+iv$ this reads $u_y+iv_y=iu_x-v_x$, that is,
$u_x=v_y$ and $u_y=-v_x$: the Cauchy--Riemann conditions
(Section 11.2) for $f$ to be holomorphic.
Under that condition it is locally exact. One can construct a primitive
on a small rectangle with sides parallel to the axes, centred at a point
$(x_0,y_0)$ and contained in the domain, without presupposing a complex
antiderivative. For $(x,y)$ in the rectangle, the horizontal segment
from $(x_0,y_0)$ to $(x,y_0)$ and the vertical segment from $(x,y_0)$
to $(x,y)$ stay inside it, so both integrals in
\[
\Phi(x,y)=\int_{x_0}^{x}f(s,y_0)\,ds
          +i\int_{y_0}^{y}f(x,t)\,dt
\]
are defined. Only the second integral depends on $y$, so the
fundamental theorem of calculus gives $\Phi_y$. For $\Phi_x$, differentiate the second
integral under the integral sign (allowed for $C^1$ $f$), then use
$if_x=f_y$ and the fundamental theorem in $t$:
\[
\Phi_y=if(x,y),\qquad
\Phi_x=f(x,y_0)+i\int_{y_0}^{y}f_x(x,t)\,dt
=f(x,y_0)+\int_{y_0}^{y}f_y(x,t)\,dt=f(x,y).
\]
Hence $d\Phi=\Phi_x\,dx+\Phi_y\,dy=f\,(dx+i\,dy)=f\,dz$. On a simply
connected complex domain, every closed curve $\Gamma$ is contractible
(it can be shrunk continuously to a point without leaving the domain).
The shrinking sweeps out a surface $S$ in the domain whose boundary is
$\Gamma$, and $d(f\,dz)=0$ on $S$, so Stokes' theorem gives
$\oint_\Gamma f\,dz=\int_Sd(f\,dz)=0$ (this is Cauchy's integral
theorem, Section 11.3). Two paths with the same endpoints together
form a closed curve, so they have equal integrals.
The local primitives thus give a global primitive
\[
\boxed{F(z)=\int_{z_0}^{z}f(\zeta)\,d\zeta+C,\qquad dF=f(z)\,dz},
\]
where the integral is path independent. To see that $dF=f\,dz$, fix a
point $z_1$ and let $\Phi$ be a local primitive, constructed as above,
on a small rectangle about $z_1$ inside the domain. For $z$ in that
rectangle, path independence lets us integrate from $z_0$ to $z_1$ and
then from $z_1$ to $z$ inside the rectangle, where $f\,d\zeta=d\Phi$;
hence
\[
F(z)=F(z_1)+\int_{z_1}^{z}d\Phi=F(z_1)+\Phi(z)-\Phi(z_1).
\]
So near $z_1$, $F$ differs from $\Phi$ by a constant, and
$dF=d\Phi=f\,dz$. More generally, this formula defines a global
primitive precisely when $\oint_\Gamma f(z)\,dz=0$ for every closed
curve $\Gamma$ in the domain. The domain condition cannot be omitted:
$f(z)=1/z$ is holomorphic on $\mathbb C\setminus\{0\}$, but on
$z=Re^{it}$, $0\le t\le2\pi$, where $dz=iRe^{it}\,dt$,
\[
\oint\frac{dz}{z}=\int_0^{2\pi}i\,dt=2\pi i\ne0,
\]
so it is not globally exact there. If $f$ is not holomorphic, even
closedness can fail: $f=\bar z=x-iy$ has $f_x=1$, $f_y=-i$, so
$d(f\,dz)=(i+i)\,dx\wedge dy=2i\,dx\wedge dy$.
\erratum{Accordingly, the Instructor's Manual's unqualified claim that
$f(z)\,dz$ is exact must be read with both the analyticity and domain
conditions above.}
